% ============================================================
%  TFM — Estudio Comparativo del Espacio Latente de Modelos
%         Generativos de Texto-Imagen
%  Autor: Víctor Menéndez Ríos
% ============================================================
\documentclass[12pt, a4paper, twoside]{book}

% --- Encoding & Language -----------------------------------
\usepackage[utf8]{inputenc}
\usepackage[T1]{fontenc}
\usepackage[spanish]{babel}

% --- Typography --------------------------------------------
\usepackage{lmodern}
\usepackage{microtype}
\usepackage{setspace}
\onehalfspacing

% --- Maths -------------------------------------------------
\usepackage{amsmath, amssymb, amsthm}
\usepackage{bm}

% --- Layout ------------------------------------------------
\usepackage[top=3cm, bottom=3cm, inner=3.5cm, outer=2.5cm]{geometry}
\usepackage{emptypage}
\pagestyle{plain}

% --- Figures & Tables --------------------------------------
\usepackage{graphicx}
\usepackage{booktabs}
\usepackage{multirow}
\usepackage{array}
\usepackage{longtable}
\usepackage{caption}
\usepackage{subcaption}
\usepackage{float}
\usepackage{xcolor}

% --- Links (sin color) -------------------------------------
\usepackage[hidelinks]{hyperref}
\usepackage{url}

% --- Listas ------------------------------------------------
\usepackage{enumitem}

% --- Custom commands ---------------------------------------
\newcommand{\z}{\mathbf{z}}
\newcommand{\x}{\mathbf{x}}
\newcommand{\eps}{\boldsymbol{\epsilon}}
\newcommand{\R}{\mathbb{R}}
\newcommand{\E}{\mathbb{E}}
\newcommand{\norm}[1]{\left\|#1\right\|_2}
\newcommand{\TODO}[1]{\textcolor{red}{\textbf{[PENDIENTE: #1]}}}
\newcommand{\PLACEHOLDER}[1]{\textcolor{orange}{\textit{[#1]}}}

% ===========================================================
\begin{document}
% ===========================================================

% --- Portada -----------------------------------------------
\begin{titlepage}
  \centering
  \vspace*{2cm}
  {\Large Universidad Politécnica de Madrid\par}
  \vspace{0.5cm}
  {\large Máster en Deep Learning\par}
  \vspace{3cm}
  \rule{\linewidth}{0.5pt}\\[0.4cm]
  {\LARGE\bfseries Estudio Comparativo del Espacio Latente\\[0.3cm]
   de Modelos Generativos de Texto-Imagen\par}
  \rule{\linewidth}{0.5pt}\\[1.5cm]
  {\large\itshape Trabajo de Fin de Máster\par}
  \vspace{2cm}
  {\large Autor: \textbf{Victor Mendez Riosalido}\par}
  \vspace{0.5cm}
  {\large Tutor: \PLACEHOLDER{Nombre del tutor}\par}
  \vfill
  {\large Curso académico 2025--2026\par}
\end{titlepage}

% --- Páginas preliminares ----------------------------------
\frontmatter

\chapter*{Resumen}
\addcontentsline{toc}{chapter}{Resumen}

Los modelos de difusión latente han transformado la síntesis de imágenes a partir de texto, pero la estructura interna de su espacio latente permanece insuficientemente comprendida. Este trabajo presenta un análisis comparativo cuantitativo del espacio latente $\z_0$ de tres modelos de la familia Stable Diffusion ---SD~1.5, SD~2.1 y SDXL--- con distintos codificadores de texto y capacidades de UNet, manteniendo constante la arquitectura del autoencoder variacional. El estudio se articula en tres ejes: análisis geométrico de las distancias latentes, análisis semántico de las direcciones de atributo y análisis del impacto del fine-tuning mediante LoRA sobre un corpus de pintura clásica (WikiArt, 500 imágenes). Los resultados muestran que SDXL presenta propiedades cualitativamente distintas en pretrain: mayor separabilidad semántica, direcciones de atributo escena-independientes con consistencia intra-tipo de 0,71 para el color, y concentración del 95\% de la energía semántica en los diez primeros componentes principales. Los tres modelos exhiben una dimensionalidad intrínseca de aproximadamente doce dimensiones efectivas, con independencia de su dimensión nominal. El fine-tuning sobre WikiArt revela un patrón contrario al esperado: SDXL, con la representación pretrain más organizada, sufre la mayor perturbación geométrica y semántica ---desplazamiento latente medio 3,5 veces superior al de los otros dos modelos, alineación PCA post-tuning de 0,54 frente a 0,82--0,89, y caídas de consistencia del 22--52\% según el atributo--- mientras que SD~1.5 y SD~2.1 muestran una deriva prácticamente nula. Este resultado evidencia que una organización semántica más estructurada no protege frente a la disrupción por adaptación de dominio, sino que la amplifica.

\bigskip
\noindent\textbf{Palabras clave:} modelos de difusión latente, espacio latente, análisis semántico, fine-tuning, LoRA, Stable Diffusion, text-to-image.

\chapter*{Abstract}
\addcontentsline{toc}{chapter}{Abstract}

Latent diffusion models have transformed text-to-image synthesis, yet the internal structure of their latent space remains poorly understood. This work presents a quantitative comparative analysis of the $\z_0$ latent space of three models from the Stable Diffusion family---SD~1.5, SD~2.1, and SDXL---differentiated by their text encoder and UNet capacity while keeping the variational autoencoder architecture constant. The study is organized around three axes: geometric analysis of latent distances, semantic analysis of attribute directions, and analysis of the impact of LoRA fine-tuning on a classical painting corpus (WikiArt, 500 images). Pre-training results show that SDXL exhibits qualitatively different properties: greater semantic separability, scene-independent attribute directions with an intra-type consistency of 0.71 for color, and concentration of 95\% of semantic energy in the top ten principal components. All three models display an intrinsic dimensionality of approximately twelve effective dimensions, regardless of their nominal dimension. WikiArt fine-tuning reveals a pattern contrary to initial predictions: SDXL, with the most organized pre-training representation, suffers the largest geometric and semantic perturbation---mean latent displacement 3.5$\times$ larger than the other two models, post-tuning PCA alignment of 0.54 versus 0.82--0.89, and consistency drops of 22--52\% across attributes---while SD~1.5 and SD~2.1 show near-negligible drift. This result demonstrates that a more structured semantic organization does not protect against disruption by domain adaptation, but amplifies it.

\bigskip
\noindent\textbf{Keywords:} latent diffusion models, latent space, semantic analysis, fine-tuning, LoRA, Stable Diffusion, text-to-image.

\tableofcontents
\listoftables
\listoffigures

% -----------------------------------------------------------
\mainmatter
% -----------------------------------------------------------

% ===========================================================
%  INTRODUCCIÓN  (incluye Estado del Arte)
% ===========================================================
\chapter{Introducción}
\label{ch:intro}

\section{Contexto y motivación}
\label{sec:motivacion}

La síntesis de imágenes a partir de descripciones textuales ha experimentado una transformación sin precedentes durante los últimos años. En agosto de 2022, la publicación de Stable Diffusion como modelo de código abierto marcó un punto de inflexión: por primera vez, un sistema capaz de producir imágenes de calidad comparable a los mejores sistemas propietarios era accesible para cualquier usuario con una tarjeta gráfica de consumo (Rombach et al., 2022). Desde entonces, el ecosistema ha crecido a un ritmo excepcional: modelos como SDXL (Podell et al., 2023), Stable Diffusion 3 (Esser et al., 2024) y Flux (Black Forest Labs, 2024) han elevado sucesivamente el techo de calidad en el ámbito del código abierto, mientras que sistemas propietarios como DALL-E 3 (Betker et al., 2023), Midjourney v6 o Adobe Firefly han redefinido los estándares de usabilidad y fidelidad visual en entornos comerciales.

Sin embargo, a pesar de la impresionante calidad de las imágenes generadas, la comprensión profunda de los mecanismos internos que sustentan la síntesis permanece notablemente incompleta. En particular, el \textit{espacio latente} ---la representación comprimida y abstracta del contenido visual que los modelos construyen durante el entrenamiento--- es la pieza central que determina la calidad, la coherencia semántica y la controlabilidad de la generación. La estructura geométrica y semántica de este espacio está lejos de ser completamente comprendida, especialmente en el contexto de los modelos de difusión latente, donde la representación latente resulta menos intuitiva que en arquitecturas anteriores como las redes generativas adversariales (Kwon et al., 2023). Esta opacidad tiene consecuencias prácticas directas: las ediciones basadas en manipulación de latentes pueden producir artefactos impredecibles, las interpolaciones no siempre son semánticamente suaves, y el comportamiento de los modelos ante el fine-tuning sobre dominios específicos es difícil de anticipar.

El presente trabajo aborda esta brecha desde una perspectiva comparativa y cuantitativa. En lugar de evaluar los modelos por la calidad visual de sus salidas ---aspecto ya extensamente tratado en la literatura---, se propone analizar y contrastar la estructura interna del espacio latente de tres modelos de difusión latente de código abierto con arquitecturas deliberadamente distintas: Stable Diffusion 1.5, Stable Diffusion 2.1 y Stable Diffusion XL. Los resultados son relevantes tanto para la comunidad investigadora ---que necesita referencias cuantitativas para comprender cómo la arquitectura afecta a la organización de las representaciones internas--- como para los profesionales que desean maximizar la controlabilidad y la predictibilidad de las ediciones.

\section{Perspectiva histórica de la síntesis visual computacional}
\label{sec:historia}

\subsection{Generación procedural y estadística (hasta 2012)}

Los primeros sistemas de síntesis visual computacional eran enteramente procedurales: algoritmos deterministas que producían imágenes a partir de reglas matemáticas explícitas. El ruido de Perlin (Perlin, 1985), introducido para la generación de texturas naturales en producción cinematográfica, es un ejemplo paradigmático. Los métodos estadísticos de síntesis de texturas de Efros y Leung (1999) demostraron que era posible sintetizar nuevas texturas a partir de un ejemplo mediante muestreo no paramétrico. El denominador común de todos estos enfoques es que la representación del mundo visual estaba codificada de forma explícita y manual, lo que limitaba drásticamente la capacidad de generalización.

\subsection{El giro neural: representaciones aprendidas (2012--2014)}

La publicación de AlexNet (Krizhevsky et al., 2012) abrió la puerta a un nuevo paradigma: en lugar de especificar qué buscar en una imagen, los sistemas aprenderían qué representaciones son útiles directamente de los datos. La transferencia de estilo neuronal (Gatys et al., 2016) demostró que las representaciones profundas contenían información visual estructurada y manipulable. El autoencoder variacional (VAE) de Kingma y Welling (2013) formalizó el primer modelo generativo profundo con un espacio latente continuo y semánticamente estructurado que podía muestrearse de forma controlada.

\subsection{La primera generación de modelos generativos profundos (2013--2017)}

Cuatro familias de modelos generativos profundos surgieron durante este período. Los \textbf{autoencoders variacionales} (Kingma y Welling, 2013) aprendían un espacio latente continuo maximizando la cota inferior de la log-verosimilitud. Las \textbf{redes generativas adversariales} (Goodfellow et al., 2014) adoptaron un juego de suma cero entre generador y discriminador; DCGAN (Radford et al., 2015) demostró que el espacio latente tenía propiedades aritméticas sorprendentes: la operación $\z_{\text{hombre con gafas}} - \z_{\text{hombre}} + \z_{\text{mujer}} \approx \z_{\text{mujer con gafas}}$ funcionaba de forma consistente. Los \textbf{modelos autorregresivos} (van den Oord et al., 2016) y los \textbf{flujos normalizantes} (Dinh et al., 2017; Kingma y Dhariwal, 2018) completaron el panorama generativo de la época.

\subsection{La era GAN: síntesis de alta fidelidad (2017--2021)}

StyleGAN (Karras et al., 2019) y StyleGAN2 (Karras et al., 2020) introdujeron la red de mapeo $f: \z \rightarrow \mathbf{w}$ que transforma el ruido latente en un espacio intermedio $\mathcal{W}$ con disentanglement superior. Trabajos como InterFaceGAN (Shen et al., 2020) y GANSpace (Härkönen et al., 2020) demostraron que atributos semánticos complejos ---edad, género, expresión, pose, iluminación--- corresponden a hiperplanos lineales o componentes principales bien definidos en $\mathcal{W}$. Esta capacidad de análisis y edición lineal en el espacio latente es el principal antecedente metodológico del presente trabajo.

\subsection{El pivote hacia la difusión y la generación guiada por texto (2021--2022)}

CLIP (Radford et al., 2021) entrenó, mediante aprendizaje contrastivo a escala masiva sobre 400 millones de pares imagen-texto, un espacio de embedding compartido entre modalidades visuales y lingüísticas. Rombach et al.\ (2022) introdujeron los \textit{modelos de difusión latente} (LDM), desplazando el proceso de difusión del espacio de píxeles al espacio latente comprimido de un VAE. Stable Diffusion, su implementación de código abierto entrenada sobre LAION-5B (Schuhmann et al., 2022), se convirtió en el primer sistema de generación texto-imagen de alta calidad accesible sin infraestructura computacional extraordinaria.

\subsection{Consolidación y madurez del ecosistema (2022--2025)}

SDXL (Podell et al., 2023) introdujo un UNet de 2.600 millones de parámetros con doble codificador de texto y resolución nativa de $1024 \times 1024$ píxeles. Stable Diffusion 3 (Esser et al., 2024) supuso un cambio de paradigma hacia transformadores de difusión multimodales y \textit{flow matching} rectificado (Lipman et al., 2022; Liu et al., 2022). Flux (Black Forest Labs, 2024), con 12.000 millones de parámetros, estableció el nuevo estado del arte en texto-imagen de código abierto a mediados de 2024. En el ámbito propietario, DALL-E 3 (Betker et al., 2023) demostró que la estrategia de \textit{recaptioning} con modelos de lenguaje mejora drásticamente la fidelidad textual sin cambios arquitectónicos.

\section{Fundamentos de los modelos generativos profundos}
\label{sec:fundamentos}

\subsection{Autoencoders variacionales}

El autoencoder variacional (Kingma y Welling, 2013) aprende simultáneamente un codificador probabilístico $q_\phi(\z \mid \x)$ y un decodificador $p_\theta(\x \mid \z)$. La función objetivo es la \textit{cota inferior de evidencia} (ELBO):

\begin{equation}
\mathcal{L}_{\text{ELBO}} = \E_{q_\phi(\z \mid \x)}\!\left[\log p_\theta(\x \mid \z)\right]
  - D_{\mathrm{KL}}\!\left[q_\phi(\z \mid \x) \,\|\, p(\z)\right]
\label{eq:elbo}
\end{equation}

La regularización mediante la divergencia KL impone una estructura global que hace el espacio latente interpolable: puntos cercanos en $\mathcal{Z}$ corresponden a imágenes perceptualmente similares.

\subsection{Redes generativas adversariales}

Las GANs (Goodfellow et al., 2014) definen un juego de suma cero entre generador $G: \z \rightarrow \x$ y discriminador $D: \x \rightarrow [0,1]$:

\begin{equation}
\min_G \max_D \;\E_{\x \sim p_{\text{data}}}\!\left[\log D(\x)\right]
+ \E_{\z \sim p(\z)}\!\left[\log\!\left(1 - D(G(\z))\right)\right]
\end{equation}

La relevancia de las GANs para este trabajo es metodológica: gran parte del instrumental de análisis del espacio latente ---búsqueda de direcciones semánticas lineales, PCA de activaciones, clasificadores lineales--- fue desarrollado inicialmente para GANs y se traslada directamente al estudio de los LDMs.

\subsection{Modelos de difusión probabilísticos}

Los modelos de difusión probabilísticos (Ho et al., 2020) definen un proceso hacia adelante que añade ruido gaussiano en $T$ pasos:

\begin{equation}
q(\x_t \mid \x_0) = \mathcal{N}\!\left(\sqrt{\bar{\alpha}_t}\,\x_0,\; (1 - \bar{\alpha}_t)\,\mathbf{I}\right),
\quad \bar{\alpha}_t = \prod_{s=1}^t (1 - \beta_s)
\end{equation}

La red neuronal $\eps_\theta$ se entrena para predecir el ruido añadido:

\begin{equation}
\mathcal{L}_{\text{simple}} = \E_{\x_0,\,\eps \sim \mathcal{N}(\mathbf{0},\mathbf{I}),\,t}\!
\left[\left\|\eps - \eps_\theta\!\left(\sqrt{\bar{\alpha}_t}\,\x_0 + \sqrt{1 - \bar{\alpha}_t}\,\eps,\;t\right)\right\|_2^2\right]
\end{equation}

El muestreo DDIM (Song et al., 2021) permite recorrer la trayectoria en $S \ll T$ pasos con naturaleza determinista, habilitando la inversión de imágenes reales. El guiado clasificador-libre (Ho y Salimans, 2022) interpola entre las predicciones de ruido condicionada e incondicionada:

\begin{equation}
\hat{\eps}_\theta(\x_t, t, y) = \eps_\theta(\x_t, t, \emptyset)
+ w\bigl[\eps_\theta(\x_t, t, y) - \eps_\theta(\x_t, t, \emptyset)\bigr]
\end{equation}

donde $w = 7{,}5$ en todos los experimentos de este trabajo.

\subsection{Modelos de difusión latente}

Rombach et al.\ (2022) propusieron desplazar el proceso de difusión al espacio latente comprimido de un VAE. Para un factor de compresión $f = 8$ y $C = 4$ canales latentes, una imagen de $512 \times 512$ se comprime en $4 \times 64 \times 64$ ---un factor de reducción dimensional de 192. El espacio latente $\z_0$ ---el latente denoised final antes de la decodificación--- es el objeto central de análisis en este trabajo. Sus propiedades geométricas y semánticas reflejan directamente cómo el modelo ha organizado el contenido visual durante el proceso de difusión.

\subsection{Transformadores de difusión y flow matching}

DiT (Peebles y Xie, 2023) propuso sustituir el UNet por un transformer que opera sobre parches del latente, siguiendo leyes de escalado predecibles sin evidencia de saturación. El \textit{flow matching} (Lipman et al., 2022; Liu et al., 2022) define trayectorias lineales entre el ruido y los datos:

\begin{equation}
\mathcal{L}_{\text{FM}} = \E_{t, \x_0, \eps}\!\left[\left\|v_\theta((1-t)\x_0 + t\eps, t) - (\eps - \x_0)\right\|_2^2\right]
\end{equation}

\section{Estado del arte en análisis del espacio latente}
\label{sec:sota}

\subsection{El espacio $\mathcal{W}$ de StyleGAN y la hipótesis de linealidad}

El análisis estructural del espacio latente encontró en StyleGAN (Karras et al., 2019) su objeto de estudio más fructífero. La hipótesis de linealidad ---que los conceptos semánticos corresponden a regiones linealmente separables--- ha demostrado ser empíricamente robusta en $\mathcal{W}$. InterFaceGAN (Shen et al., 2020) entrenó una SVM lineal para encontrar el hiperplano que separa imágenes con y sin un atributo dado:

\begin{equation}
\langle \mathbf{n}, \mathbf{w} \rangle + b = 0
\end{equation}

La edición se realiza como $\mathbf{w}' = \mathbf{w} + \alpha\mathbf{n}$. La contribución más relevante es que estas direcciones son transferibles: vectores encontrados sobre un subconjunto de latentes se aplican con alta calidad a latentes arbitrarios, confirmando que corresponden a propiedades globales de la distribución aprendida.

GANSpace (Härkönen et al., 2020) demostró que el PCA sobre vectores $\mathbf{w}$ produce componentes interpretables de forma no supervisada. SeFa (Shen y Zhou, 2021) mostró que los vectores propios de la primera capa del generador son suficientes para recuperar estas direcciones sin necesidad de datos. Wu et al.\ (2021) propusieron el \textit{StyleSpace} $\mathcal{S}$ formado por los coeficientes de modulación canal a canal, con un disentanglement superior al de $\mathcal{W}$.

\subsection{El espacio latente en los modelos de difusión latente}

La inversión DDIM (Song et al., 2021) proporciona la base técnica para recuperar el latente que genera una imagen dada; Mokady et al.\ (2023) refinaron este procedimiento con la \textit{inversión de texto nulo} para garantizar reconstrucciones exactas. Kwon et al.\ (2023) identificaron el \textit{h-space} ---activaciones del bottleneck del UNet--- como un espacio con linealidad semántica análoga al $\mathcal{W}$ de StyleGAN:

\begin{equation}
\mathbf{h}_t = \text{UNet}_\text{enc}(\z_t, t, \tau_\theta(y))
\end{equation}

La operación de edición $\mathbf{h}'_t = \mathbf{h}_t + \alpha\Delta\mathbf{h}$ produce cambios semánticos coherentes que transfieren a imágenes no vistas. Sin embargo, el trabajo se limita a un único modelo sin estudiar la dependencia con la arquitectura. Hertz et al.\ (2022) demostraron que los mapas de atención cruzada codifican la disposición espacial de los conceptos, habilitando Prompt-to-Prompt. Tang et al.\ (2023) identificaron que las activaciones del UNet en pasos de tiempo intermedios superan a DINO y CLIP para correspondencia semántica densa.

\subsection{Dimensionalidad intrínseca de los espacios latentes}

Pope et al.\ (2021) estudiaron la dimensionalidad intrínseca de los espacios latentes usando el estimador de dos vecinos más cercanos (TwoNN, Facco et al., 2017), encontrando que el espacio $\mathcal{W}$ de StyleGAN para rostros tiene una dimensionalidad intrínseca de aproximadamente 40--60 dimensiones efectivas pese a sus 512 nominales. Esta propiedad de baja dimensionalidad es una hipótesis verificable también para los LDMs, cuya dimensión nominal oscila entre 16.384 y 65.536 elementos.

\subsection{Fine-tuning y personalización}

La inversión textual (Gal et al., 2022) optimiza el embedding de un token virtual sin modificar los pesos del modelo. DreamBooth (Ruiz et al., 2023) fine-tunea el UNet completo con pérdida de preservación de clases para evitar el olvido catastrófico (McCloskey y Cohen, 1989; Kirkpatrick et al., 2017). LoRA (Hu et al., 2022) factoriza el desplazamiento de pesos en bajo rango:

\begin{equation}
\mathbf{W} = \mathbf{W}_0 + \mathbf{B}\mathbf{A},
\quad \mathbf{B} \in \R^{d \times r},\; \mathbf{A} \in \R^{r \times k},\; r \ll \min(d, k)
\end{equation}

El análisis de cómo el fine-tuning modifica la estructura interna del espacio latente es el aspecto de mayor novedad de este trabajo respecto a la literatura existente.

\subsection{Brechas identificadas y posicionamiento}

La revisión permite identificar cuatro brechas principales. Primera: la práctica totalidad de los estudios de análisis del espacio latente en LDMs se realizan sobre SD~1.x, sin comparación sistemática entre arquitecturas. Segunda: la literatura se centra en el h-space y los mapas de atención, no en el espacio $\z_0$ producido por el VAE, más accesible para aplicaciones de edición basadas en inversión DDIM. Tercera: los métodos de personalización se han evaluado en términos de calidad visual, pero no en términos de las propiedades estructurales del espacio latente resultante. Cuarta: no existe un protocolo replicable para la evaluación comparativa del espacio latente de LDMs aplicable uniformemente a distintas arquitecturas.

\section{Objetivos}
\label{sec:objetivos}

El objetivo general es caracterizar y comparar cuantitativamente la estructura del espacio latente $\z_0$ de SD~1.5, SD~2.1 y SDXL, y analizar cómo dicha estructura se modifica tras el fine-tuning sobre un dominio visual específico. Los objetivos específicos son: (1) cuantificar la relación entre distancia latente y similitud perceptual y caracterizar la dimensionalidad intrínseca; (2) identificar y cuantificar direcciones semánticas para atributos visuales controlados y evaluar el grado de disentanglement; (3) medir la deriva geométrica y semántica del espacio latente tras LoRA sobre WikiArt y comparar la respuesta de los tres modelos; y (4) identificar qué propiedades del espacio latente son invariantes a la arquitectura y cuáles dependen del codificador de texto, la capacidad del UNet o la resolución nativa.

% ===========================================================
%  METODOLOGÍA
% ===========================================================
\chapter{Metodología}
\label{ch:metodos}

\section{Entorno experimental}
\label{sec:entorno}

\subsection{Plataforma de cómputo}

Los experimentos se ejecutaron sobre Google Colaboratory (Google LLC, 2017), utilizando la GPU NVIDIA Tesla T4 (15 GB VRAM, arquitectura Turing). La decisión de emplear computación en la nube responde a las limitaciones del equipo local (NVIDIA GeForce GTX~1050~Max-Q, 3~GB~VRAM), insuficiente para ejecutar los modelos de difusión en condiciones reproducibles. La persistencia de datos entre sesiones se gestiona mediante Google Drive. El equipo local (Intel Core i5-1035G1, 7,5~GB~RAM, Ubuntu~22.04) se utiliza exclusivamente para el desarrollo del código.

\subsection{Entorno de software y reproducibilidad}

La Tabla~\ref{tab:software} recoge las versiones exactas de las dependencias principales. La reproducibilidad se garantiza mediante especificación explícita de versiones en \texttt{requirements\_colab.txt}, un cuaderno de entrada estandarizado que instala las dependencias al inicio de cada sesión, control de versiones con Git, y semillas aleatorias fijas en PyTorch, NumPy y Python.

\begin{table}[htb]
\centering
\caption{Versiones exactas de las dependencias del entorno experimental.}
\label{tab:software}
\small
\begin{tabular}{lll}
\toprule
\textbf{Biblioteca} & \textbf{Versión} & \textbf{Propósito} \\
\midrule
Python       & 3.11.15       & Lenguaje base \\
PyTorch      & 2.3.1+cu121   & Framework de aprendizaje profundo \\
Diffusers    & 0.27.2        & Pipelines de modelos de difusión \\
Transformers & 4.40.2        & Codificadores de texto \\
Accelerate   & 0.30.1        & Entrenamiento multi-dispositivo \\
PEFT         & 0.11.1        & Fine-tuning eficiente (LoRA) \\
xformers     & 0.0.27        & Atención eficiente en memoria \\
open-clip-torch & 2.24.0    & Modelo CLIP de código abierto \\
scikit-learn & 1.8.0         & Análisis estadístico y PCA \\
\bottomrule
\end{tabular}
\end{table}

\subsection{Métricas de evaluación}

Se utilizan dos métricas de similitud de imagen. El MSE mide la diferencia cuadrática media píxel a píxel y sirve como referencia global de disimilitud. El SSIM global (Wang et al., 2004) evalúa la similitud estructural entre las estadísticas de las imágenes y constituye la métrica perceptual principal del análisis de correlación latente--imagen.

\section{Modelos bajo estudio}
\label{sec:modelos}

\subsection{Marco teórico: modelos de difusión latente}

Los modelos de difusión latente (Rombach et al., 2022) operan en dos etapas: compresión perceptual mediante un VAE y difusión en el espacio latente comprimido. El proceso de difusión en el espacio latente sigue:

\begin{equation}
q(\z_t \mid \z_0) = \mathcal{N}\!\left(\sqrt{\bar{\alpha}_t}\,\z_0,\; (1 - \bar{\alpha}_t)\,\mathbf{I}\right)
\end{equation}

El modelo de ruido $\eps_\theta$ se entrena minimizando:

\begin{equation}
\mathcal{L} = \E_{\z_0,\, \eps \sim \mathcal{N}(0,\mathbf{I}),\, t}\!
\left[\norm{\eps - \eps_\theta(\z_t, t, \tau_\theta(y))}^2\right]
\end{equation}

donde $\tau_\theta(y)$ es la representación del prompt producida por el codificador de texto.

\subsection{Criterios de selección}

La selección se guió por tres criterios. \textbf{Diversidad arquitectónica}: los modelos difieren en el codificador de texto y en la capacidad del UNet, las dos variables que se desean estudiar. \textbf{VAE compartido}: los tres modelos utilizan la misma arquitectura de autoencoder variacional con factor de compresión $f = 8$ y $C = 4$ canales latentes, garantizando que las diferencias observadas en el espacio $\z_0$ son atribuibles al proceso de difusión y al condicionamiento textual. \textbf{Viabilidad computacional} dentro de los 15~GB~VRAM de la GPU~T4.

\subsection{Stable Diffusion 1.5}

SD~1.5 (Rombach et al., 2022) fue entrenado por RunwayML sobre LAION-Aesthetics v2 5+. El codificador de texto es CLIP ViT-L/14 (Radford et al., 2021), que proyecta el prompt en 77 tokens de 768 dimensiones. El UNet tiene aproximadamente 860 millones de parámetros y la resolución nativa es $512 \times 512$ con latentes de forma $4 \times 64 \times 64$.

\subsection{Stable Diffusion 2.1}

SD~2.1 (Stability AI, 2022) introduce dos diferencias respecto a SD~1.5. El codificador OpenCLIP ViT-H/14 (Ilharco et al., 2021) produce embeddings de 1.024 dimensiones frente a las 768 de CLIP ViT-L/14. La resolución nativa es $768 \times 768$ con latentes de forma $4 \times 96 \times 96$.

\subsection{Stable Diffusion XL}

SDXL (Podell et al., 2023) introduce tres cambios. El \textbf{doble codificador de texto} ---CLIP ViT-L/14 más OpenCLIP ViT-bigG/14 (1.300 millones de parámetros)--- produce un vector de condicionamiento de 2.048 dimensiones de atención cruzada más 1.280 de embedding global. El \textbf{UNet de 2.600 millones de parámetros} es $\sim 3\times$ mayor que el de los modelos anteriores. La \textbf{resolución nativa de $1024 \times 1024$} produce latentes de forma $4 \times 128 \times 128$.

\begin{table}[htb]
\centering
\caption{Comparativa arquitectónica de los modelos bajo estudio.}
\label{tab:arquitecturas}
\small
\begin{tabular}{lccc}
\toprule
\textbf{Propiedad} & \textbf{SD 1.5} & \textbf{SD 2.1} & \textbf{SDXL} \\
\midrule
Codificador de texto & CLIP ViT-L/14 & OpenCLIP ViT-H/14 & CLIP ViT-L + OpenCLIP ViT-bigG \\
Dim.\ embedding texto & 768 & 1.024 & 2.048 + 1.280 \\
Parámetros UNet (aprox.) & 860 M & 870 M & 2.600 M \\
Resolución nativa & $512\times512$ & $768\times768$ & $1024\times1024$ \\
Factor de compresión ($f$) & 8 & 8 & 8 \\
Canales latentes ($C$) & 4 & 4 & 4 \\
Forma del latente & $4\times64\times64$ & $4\times96\times96$ & $4\times128\times128$ \\
Dimensión latente total & 16.384 & 36.864 & 65.536 \\
\bottomrule
\end{tabular}
\end{table}

\subsection{Interfaz unificada y extracción de $\z_0$}

Se implementó la clase \texttt{ModelWrapper} que encapsula el pipeline de Diffusers y expone tres operaciones estables: codificación de imagen a latente, decodificación de latente a imagen, y generación con captura del latente final. El latente $\z_0$ en modo codificación se extrae como la media de la distribución posterior multiplicada por el factor de escalado del VAE, garantizando determinismo. En modo generación se usa un hook temporal sobre \texttt{vae.decode} para capturar $\z_0$ antes de la decodificación.

\section{Corpus de vectores latentes}
\label{sec:corpus}

\subsection{Diseño}

Se generaron 150 vectores latentes por modelo, organizados en dos categorías: \textbf{latentes base} (60 por modelo, 20 prompts $\times$ 3 semillas) para el análisis geométrico global y la base PCA del análisis semántico; y \textbf{latentes de atributo} (90 por modelo, 15 grupos $\times$ 2 valores $\times$ 3 semillas) para construir las direcciones semánticas. Los mismos prompts, semillas e hiperparámetros se aplican a los tres modelos.

\subsection{Diseño del banco de prompts}

Los 20 prompts base y los 30 prompts de atributo se diseñaron con el asistente de inteligencia artificial Claude (Anthropic, 2024), empleándolo para generar descripciones diversas, equilibradas y libres de correlación temática inadvertida. Los prompts base cubren una amplia diversidad de dominios visuales (naturaleza muerta, fauna, arquitectura, gastronomía, paisaje, espacio urbano, tecnología), escalas compositivas y condiciones de luz. La lista completa se recoge en el Apéndice~\ref{app:prompts}.

Los 15 grupos de atributo se construyeron como \textbf{pares controlados}: ambos prompts del mismo grupo describen la misma escena modificando exclusivamente el atributo objetivo. Los tres tipos de atributo son color (temperatura cálida frente a fría, 5 grupos), iluminación (alta clave frente a baja clave, 5 grupos) y estilo (fotografía fotorrealista frente a imagen artística, 5 grupos). La dirección semántica de un grupo $g$ se define como:

\begin{equation}
\mathbf{d}_g = \bar{\z}(g, v^+) - \bar{\z}(g, v^-)
\label{eq:dir_semantica}
\end{equation}

donde $\bar{\z}(g, v)$ es la media sobre las tres semillas, eliminando la variabilidad estocástica para aislar el efecto del atributo.

\subsection{Hiperparámetros de generación}

\begin{table}[htb]
\centering
\caption{Hiperparámetros de generación del corpus.}
\label{tab:hiper_generacion}
\small
\begin{tabular}{lccc}
\toprule
\textbf{Parámetro} & \textbf{SD 1.5} & \textbf{SD 2.1} & \textbf{SDXL} \\
\midrule
Sampler & DDIM & DDIM & DDIM \\
Pasos de denoising & 50 & 50 & 30 \\
Guidance scale $w$ & 7,5 & 7,5 & 7,5 \\
Semillas & 0, 1, 2 & 0, 1, 2 & 0, 1, 2 \\
Resolución de salida & $512^2$ & $512^2$ & $512^2$ \\
Precisión & fp16 & fp16 & fp16 (VAE float32) \\
\bottomrule
\end{tabular}
\end{table}

\section{Análisis geométrico}
\label{sec:metodos:geometrico}

\subsection{Separabilidad intra/inter-prompt}

Para cada modelo se construye la matriz de distancias euclídeas $\mathbf{D} \in \R^{N \times N}$ sobre los $N = 60$ latentes base. El \textbf{ratio de separabilidad} $\text{sep} = \mu_{\text{inter}} / \mu_{\text{intra}}$ mide si los latentes del mismo prompt están más agrupados que los de prompts distintos. La significación estadística se evalúa con el test de Kolmogorov--Smirnov bilateral.

\subsection{Dimensionalidad efectiva}

Los valores propios de la covarianza se obtienen mediante SVD thin de la matriz de datos centrada. El \textbf{participation ratio}:

\begin{equation}
\text{PR} = \frac{\left(\sum_k \lambda_k\right)^2}{\sum_k \lambda_k^2}
\end{equation}

mide el número efectivo de dimensiones activas de forma continua y sin necesidad de elegir un umbral.

\subsection{Correlación distancia latente--similitud imagen}

Se miden la correlación de Pearson $r$ y de Spearman $\rho$ entre la distancia latente L2 y las métricas de disimilitud MSE y $-\text{SSIM}$.

\section{Análisis semántico}
\label{sec:metodos:semantico}

\subsection{Análisis de componentes principales}

El PCA se aplica sobre los 60 latentes base aplanados a vectores de dimensión $D$. Las métricas extraídas son la varianza explicada por componente y el número de componentes para el 80\% y 90\% de varianza acumulada.

\subsection{Consistencia y ortogonalidad de las direcciones semánticas}

La \textbf{consistencia intra-tipo} mide si la dirección de un atributo es universal (escena-independiente):

\begin{equation}
\text{consistencia}(t) = \frac{1}{\binom{|G_t|}{2}} \sum_{g \neq g' \in G_t} \cos\!\left(\hat{\mathbf{d}}_g, \hat{\mathbf{d}}_{g'}\right)
\end{equation}

El \textbf{entrelazamiento inter-tipo} mide la ausencia de disentanglement:

\begin{equation}
\text{entrelazamiento}(t_1, t_2) = \frac{1}{|G_{t_1}| \cdot |G_{t_2}|}
\sum_{g,\, g'} \left|\cos\!\left(\hat{\mathbf{d}}_g, \hat{\mathbf{d}}_{g'}\right)\right|
\end{equation}

La fracción de energía semántica en los $K_0 = 10$ primeros componentes PCA cuantifica si el atributo se codifica en el subespacio de alta varianza:

\begin{equation}
f_g = \frac{\sum_{k=1}^{K_0} (\mathbf{V}_k^\top \hat{\mathbf{d}}_g)^2}{\sum_{k=1}^{K} (\mathbf{V}_k^\top \hat{\mathbf{d}}_g)^2}
\end{equation}

\section{Fine-tuning y análisis de deriva latente}
\label{sec:metodos:finetune}

\subsection{Diseño experimental}

El análisis de fine-tuning cierra el ciclo del trabajo: los espacios latentes caracterizados en las secciones anteriores se someten a adaptación de dominio y se vuelven a caracterizar con los mismos instrumentos. Los tres modelos se adaptan sobre el mismo dataset, con los mismos hiperparámetros y el mismo número de pasos de actualización.

\subsection{Selección del dominio: WikiArt}

Se selecciona el dominio de \textbf{pintura clásica} utilizando WikiArt (Saleh y Elgammal, 2015) por cuatro razones. Primero, ataca directamente el atributo con consistencia intra-tipo más baja en los tres modelos (estilo: 0,028 en SD~1.5, 0,010 en SD~2.1, 0,209 en SDXL), permitiendo verificar si la exposición intensiva a un dominio artístico induce la aparición de una dirección de estilo más universal. Segundo, la distribución de pinturas artísticas es cualitativamente diferente de la fotografía dominante en LAION-5B (Schuhmann et al., 2022), maximizando la señal detectable en las métricas de deriva. Tercero, dado que SDXL ya tiene consistencia de estilo un orden de magnitud mayor que SD~1.5 y SD~2.1, existe una hipótesis diferencial comprobable: \textit{SDXL absorberá el conocimiento artístico de WikiArt de forma más eficiente y producirá mayor amplificación de la consistencia semántica de estilo.} Cuarto, WikiArt es un dataset de referencia en la literatura de síntesis de estilos artísticos (Gatys et al., 2016).

\subsection{Dataset WikiArt}

Se selecciona un subconjunto balanceado de 500 imágenes: 50 por cada uno de los diez movimientos artísticos representados (impresionismo, realismo, romanticismo, expresionismo, post-impresionismo, barroco, cubismo, arte abstracto, simbolismo y Art Nouveau). Las captions se generan a partir de los metadatos con la plantilla \textit{``a \{style\} painting\{, \{genre\},\} by \{artist\}''}.

\subsection{Configuración LoRA}

LoRA se aplica exclusivamente al UNet; el codificador de texto y el VAE se mantienen \textbf{congelados}. La congelación del codificador de texto asegura que los cambios en los latentes sean atribuibles a la reorganización del proceso de denoising. La congelación del VAE garantiza que $\z_0^\text{pre}$ y $\z_0^\text{ft}$ sean comparables en el mismo espacio vectorial.

\begin{table}[htb]
\centering
\caption{Configuración LoRA y hiperparámetros de entrenamiento.}
\label{tab:lora}
\small
\begin{tabular}{lccc}
\toprule
\textbf{Parámetro} & \textbf{SD 1.5} & \textbf{SD 2.1} & \textbf{SDXL} \\
\midrule
Rango ($r$) & 8 & 8 & 8 \\
Alpha ($\alpha$) & 8 & 8 & 8 \\
Dropout & 0,1 & 0,1 & 0,1 \\
Módulos objetivo & \multicolumn{3}{c}{\texttt{to\_q, to\_k, to\_v, to\_out.0}} \\
Parámetros entrenables & $\approx 3{,}4$ M & $\approx 3{,}5$ M & $\approx 11{,}2$ M \\
Fracción del total & 0,4\% & 0,4\% & 0,4\% \\
Pasos de entrenamiento & 1.000 & 1.000 & 800 \\
Batch efectivo & 8 & 8 & 4 \\
Learning rate & $10^{-4}$ & $10^{-4}$ & $10^{-4}$ \\
LR scheduler & cosine & cosine & cosine \\
Optimizador & AdamW & AdamW & AdamW \\
\bottomrule
\end{tabular}
\end{table}

\subsection{Métricas de deriva latente}

Tras el fine-tuning se regenera el corpus completo con el modelo adaptado. Las métricas de deriva son: $\Delta\text{sep}$, $\Delta\text{PR}$ y drift del centroide para la \textbf{deriva geométrica}; alineación de Grassmann entre los $K_0 = 10$ subespacios principales para la \textbf{rotación PCA}; y $\Delta\norm{\mathbf{d}_g}$, $\Delta\text{consistencia}$, $\Delta\text{entrelazamiento}$ y $\Delta f_g$ para la \textbf{deriva semántica}.

% ===========================================================
%  RESULTADOS
% ===========================================================
\chapter{Resultados}
\label{ch:resultados}

\section{Verificación de modelos y análisis del espacio latente $\z_0$}
\label{sec:res:modelos}

\subsection{Estadísticas descriptivas}

Los tres modelos se cargaron correctamente en la GPU T4 en precisión fp16, con la excepción del VAE de SDXL, que requirió float32 para evitar desbordamiento numérico. La Tabla~\ref{tab:estadisticas_z0} recoge las estadísticas descriptivas del latente $\z_0$ para el mismo prompt y semilla.

\begin{table}[htb]
\centering
\caption{Estadísticas descriptivas del espacio latente $\z_0$ (mismo prompt y semilla en los tres modelos).}
\label{tab:estadisticas_z0}
\small
\begin{tabular}{lrrrrr}
\toprule
\textbf{Modelo} & \textbf{Dim.} & $\bar{z}$ & $\sigma_z$ & $\norm{\z}$ & \textbf{MAE VAE} \\
\midrule
SD 1.5 & 16.384 & $-0{,}083$ & $0{,}856$ & $110{,}0$ & $4{,}00$ \\
SD 2.1 & 36.864 & $-0{,}162$ & $0{,}669$ & $132{,}2$ & $2{,}64$ \\
SDXL   & 65.536 & $-0{,}188$ & $0{,}595$ & $159{,}7$ & $1{,}80$ \\
\bottomrule
\end{tabular}
\end{table}

La norma L2 crece monótonamente aunque la desviación típica decrece, coherente con la expresión $\E[\norm{\z}] \approx \sqrt{N(\mu^2 + \sigma^2)}$, confirmando distribuciones aproximadamente gaussianas. La mejora progresiva del MAE del ciclo encode$\to$decode (4,00 $\to$ 2,64 $\to$ 1,80) indica que los modelos más recientes poseen VAEs de mayor capacidad; el error del VAE de SDXL es subperceptual (0,71\% relativo).

\subsection{Estructura espacial de los canales latentes}

\begin{figure}[htb]
\centering
\includegraphics[width=\textwidth]{data/fase1/fase1_latent_channels}
\caption{Visualización de los cuatro canales del tensor $\z_0$ para SD~1.5 (izquierda), SD~2.1 (centro) y SDXL (derecha) con el mismo prompt y semilla.}
\label{fig:canales}
\end{figure}

La inspección visual (Figura~\ref{fig:canales}) revela diferencias cualitativas entre modelos. En SD~1.5, los cuatro canales muestran la silueta del objeto principal con complementariedad sistemática entre pares. En SD~2.1, la estructura semántica se conserva con mayor rango de valores ($\approx[-10, +7{,}5]$ frente a $[-6, +4]$), reflejo del condicionamiento más expresivo del codificador de mayor dimensión. En SDXL, destaca la presencia de bandas verticales periódicas que corresponden directamente al patrón estructural del fondo de la imagen generada: este VAE codifica los patrones espaciales del fondo con una fidelidad comparable al objeto principal, evidenciando una representación más distribuida coherente con su menor MAE.

\section{Corpus de latentes generado}
\label{sec:res:corpus}

Se generaron 450 vectores latentes en total (150 por modelo), almacenados como tensores PyTorch float32 en formato $[N, 1, C, H, W]$. La verificación de completitud y reproducibilidad confirmó que la regeneración de un subconjunto de muestras con los mismos hiperparámetros produce vectores idénticos bit a bit.

\section{Análisis geométrico del espacio latente}
\label{sec:res:geometrico}

\subsection{Separabilidad intra/inter-prompt}

\begin{table}[htb]
\centering
\caption{Estadísticas de distancias L2 intra e inter-prompt (60 latentes base por modelo).}
\label{tab:separabilidad}
\small
\begin{tabular}{lrrrrrrr}
\toprule
\textbf{Modelo} & $\mu_\text{intra}$ & $\sigma_\text{intra}$ & $\mu_\text{inter}$ & $\sigma_\text{inter}$ & \textbf{sep} & \textbf{KS} & \textbf{$p$} \\
\midrule
SD 1.5 & 151,8 & 16,4 & 148,1 & 22,4 & 0,976 & 0,142 & 0,173 \\
SD 2.1 & 195,8 & 22,4 & 189,1 & 30,7 & 0,966 & 0,187 & 0,031 \\
SDXL   & 230,5 & 35,0 & 257,6 & 42,9 & \textbf{1,118} & 0,304 & $<{0{,}001}$ \\
\bottomrule
\end{tabular}
\end{table}

\begin{figure}[htbp]
\centering
\begin{subfigure}[b]{\textwidth}
    \includegraphics[width=\textwidth]{data/fase3/results/geo_sd15_distancias}
    \caption*{SD~1.5}
\end{subfigure}\\[0.4em]
\begin{subfigure}[b]{\textwidth}
    \includegraphics[width=\textwidth]{data/fase3/results/geo_sd21_distancias}
    \caption*{SD~2.1}
\end{subfigure}\\[0.4em]
\begin{subfigure}[b]{\textwidth}
    \includegraphics[width=\textwidth]{data/fase3/results/geo_sdxl_distancias}
    \caption*{SDXL}
\end{subfigure}
\caption{Distribuciones de distancias L2 intra-prompt e inter-prompt en el espacio latente de los tres modelos.}
\label{fig:distancias}
\end{figure}

SD~1.5 y SD~2.1 presentan un ratio de separabilidad inferior a 1 (0,976 y 0,966), indicando que los latentes del mismo prompt con distintas semillas están en promedio más alejados que latentes de prompts distintos (Tabla~\ref{tab:separabilidad}, Figura~\ref{fig:distancias}). El test KS no rechaza la igualdad de distribuciones para SD~1.5 ($p = 0{,}173$) y solo marginalmente para SD~2.1 ($p = 0{,}031$). En estos dos modelos, la variabilidad estocástica del proceso de denoising desplaza el latente tanto como un cambio de contenido semántico. SDXL, en cambio, muestra sep = 1,118 con $p < 0{,}001$: los latentes del mismo prompt están sistemáticamente más próximos que los de prompts distintos.

\subsection{Dimensionalidad efectiva}

\begin{table}[htb]
\centering
\caption{Dimensionalidad efectiva del espacio latente base.}
\label{tab:dimensionalidad}
\small
\begin{tabular}{lccccc}
\toprule
\textbf{Modelo} & \textbf{Dim.\ nominal} & \textbf{PR} & \textbf{dims@80\%} & \textbf{dims@90\%} & \textbf{dims@95\%} \\
\midrule
SD 1.5 & 16.384 & 12,35 & 25 & 38 & 46 \\
SD 2.1 & 36.864 & 12,33 & 27 & 39 & 47 \\
SDXL   & 65.536 & 11,56 & 22 & 35 & 45 \\
\bottomrule
\end{tabular}
\end{table}

\begin{figure}[htb]
\centering
\includegraphics[width=0.85\textwidth]{data/fase3/results/snapshot_cross_model_pca}
\caption{Varianza acumulada en función del número de componentes principales.}
\label{fig:pca_varianza}
\end{figure}

El participation ratio $\approx 12$ es prácticamente idéntico en los tres modelos pese a que sus dimensiones nominales difieren en un factor de cuatro (Tabla~\ref{tab:dimensionalidad}, Figura~\ref{fig:pca_varianza}). El número de dimensiones necesarias para capturar el 80\% de la varianza (22--27) supone menos del 0,04\% de la dimensión nominal de SDXL, confirmando que la baja dimensionalidad intrínseca es una propiedad arquitectónica compartida.

\subsection{Correlación entre distancia latente y similitud de imagen}

\begin{table}[htb]
\centering
\caption{Correlación entre distancia latente L2 y métricas de disimilitud de imagen (500 pares muestreados aleatoriamente por modelo, semilla fija 42, de los 1.770 pares posibles $\binom{60}{2}$).}
\label{tab:correlacion}
\small
\begin{tabular}{lrrrr}
\toprule
\textbf{Modelo} & $r(\text{L2, MSE})$ & $\rho(\text{L2, MSE})$ & $r(\text{L2}, -\text{SSIM})$ & \textbf{$p$-valor} \\
\midrule
SD 1.5 & 0,740 & 0,763 & 0,796 & $\ll 10^{-80}$ \\
SD 2.1 & 0,718 & 0,753 & 0,761 & $\ll 10^{-80}$ \\
SDXL   & 0,680 & 0,698 & 0,691 & $\ll 10^{-68}$ \\
\bottomrule
\end{tabular}
\end{table}

La correlación es positiva, fuerte y altamente significativa en los tres modelos (Tabla~\ref{tab:correlacion}). El orden $r_{\text{sd15}} > r_{\text{sd21}} > r_{\text{sdxl}}$ sugiere que el VAE de menor resolución de SD~1.5 actúa como un filtro de información de alta frecuencia más agresivo, produciendo latentes más correlacionados con métricas globales como el MSE.

\section{Análisis semántico del espacio latente}
\label{sec:res:semantico}

\subsection{Estructura dimensional: análisis PCA}

\begin{table}[htb]
\centering
\caption{Varianza explicada por los primeros componentes principales.}
\label{tab:pca}
\small
\begin{tabular}{lrrrrr}
\toprule
\textbf{Modelo} & \textbf{PC1 (\%)} & \textbf{PC2 (\%)} & \textbf{PC1+PC2 (\%)} & \textbf{dims@80\%} & \textbf{dims@90\%} \\
\midrule
SD 1.5 & 19,2 & 17,7 & 36,9 & 25 & 38 \\
SD 2.1 & 20,1 & 17,1 & 37,3 & 27 & 39 \\
SDXL   & 19,1 & 14,6 & 33,8 & 22 & 35 \\
\bottomrule
\end{tabular}
\end{table}

Los tres modelos comparten un rasgo característico: las dos primeras componentes concentran el 33--37\% de la varianza total. La diferencia más notable se observa en SDXL: su tercera componente explica el 13,3\% (acumulado del 47,1\% con tres componentes), frente al 6,1\% y 5,3\% de los modelos anteriores, y coincide con la dirección de iluminación.

\subsection{Discriminabilidad de las direcciones semánticas}

\begin{table}[htb]
\centering
\caption{Norma media ($\pm$ desviación típica) de las direcciones semánticas por tipo y modelo.}
\label{tab:normas_dir}
\small
\begin{tabular}{lrrr}
\toprule
\textbf{Tipo de atributo} & \textbf{SD 1.5} & \textbf{SD 2.1} & \textbf{SDXL} \\
\midrule
Color (cálido/frío)        & $65{,}4 \pm 4{,}1$  & $90{,}4 \pm 10{,}9$ & $183{,}4 \pm 18{,}8$ \\
Iluminación (brillante/oscuro) & $55{,}8 \pm 11{,}4$ & $76{,}1 \pm 13{,}5$ & $149{,}8 \pm 29{,}1$ \\
Estilo (foto/artístico)    & $51{,}6 \pm 15{,}0$ & $64{,}8 \pm 12{,}3$ & $106{,}5 \pm 32{,}7$ \\
\bottomrule
\end{tabular}
\end{table}

\begin{figure}[htbp]
\centering
\begin{subfigure}[b]{\textwidth}
    \includegraphics[width=\textwidth]{data/fase4/results/sem_sd15_direcciones}
    \caption*{SD~1.5}
\end{subfigure}\\[0.4em]
\begin{subfigure}[b]{\textwidth}
    \includegraphics[width=\textwidth]{data/fase4/results/sem_sd21_direcciones}
    \caption*{SD~2.1}
\end{subfigure}\\[0.4em]
\begin{subfigure}[b]{\textwidth}
    \includegraphics[width=\textwidth]{data/fase4/results/sem_sdxl_direcciones}
    \caption*{SDXL}
\end{subfigure}
\caption{Norma media de las direcciones semánticas por tipo de atributo y modelo.}
\label{fig:normas}
\end{figure}

Se observan dos patrones consistentes (Tabla~\ref{tab:normas_dir}, Figura~\ref{fig:normas}). La \textbf{jerarquía interna} color $>$ iluminación $>$ estilo es universal en los tres modelos. Las normas de SDXL son 2,1--2,8$\times$ las de SD~1.5, un crecimiento proporcionalmente mayor que el de las normas base (1,8$\times$), indicando mayor expresividad semántica.

\subsection{Consistencia intra-tipo de las direcciones semánticas}

\begin{table}[htb]
\centering
\caption{Similitud coseno media intra-tipo. Valores en $[-1, 1]$; mayor indica mayor universalidad de la dirección.}
\label{tab:consistencia}
\small
\begin{tabular}{lrrr}
\toprule
\textbf{Tipo} & \textbf{SD 1.5} & \textbf{SD 2.1} & \textbf{SDXL} \\
\midrule
Color       & $0{,}300 \pm 0{,}061$ & $0{,}229 \pm 0{,}112$ & $\mathbf{0{,}707 \pm 0{,}048}$ \\
Iluminación & $0{,}109 \pm 0{,}091$ & $0{,}155 \pm 0{,}062$ & $\mathbf{0{,}517 \pm 0{,}111}$ \\
Estilo      & $0{,}028 \pm 0{,}081$ & $0{,}010 \pm 0{,}082$ & $\mathbf{0{,}209 \pm 0{,}081}$ \\
\bottomrule
\end{tabular}
\end{table}

\begin{figure}[htbp]
\centering
\begin{subfigure}[b]{\textwidth}
    \includegraphics[width=\textwidth]{data/fase4/results/sem_sd15_angulos}
    \caption*{SD~1.5}
\end{subfigure}\\[0.4em]
\begin{subfigure}[b]{\textwidth}
    \includegraphics[width=\textwidth]{data/fase4/results/sem_sd21_angulos}
    \caption*{SD~2.1}
\end{subfigure}\\[0.4em]
\begin{subfigure}[b]{\textwidth}
    \includegraphics[width=\textwidth]{data/fase4/results/sem_sdxl_angulos}
    \caption*{SDXL}
\end{subfigure}
\caption{Mapa de calor de similitudes coseno entre las direcciones semánticas para SD~1.5 (izquierda), SD~2.1 (centro) y SDXL (derecha).}
\label{fig:consistencia}
\end{figure}

SD~1.5 y SD~2.1 muestran consistencias bajas para todos los tipos (Tabla~\ref{tab:consistencia}, Figura~\ref{fig:consistencia}): las direcciones semánticas no son universales sino dependientes del contexto. SDXL muestra consistencias notablemente más altas: 0,707 para color (ángulo de ${\sim}45^{\circ}$), 0,517 para iluminación y 0,209 para estilo. En SDXL existe una \textit{dirección universal del color} escena-independiente.

\subsection{Ortogonalidad inter-tipo}

\begin{table}[htb]
\centering
\caption{Similitud coseno absoluta media inter-tipo. Valores en $[0, 1]$; menor indica mejor disentanglement.}
\label{tab:ortogonalidad}
\small
\begin{tabular}{lrrr}
\toprule
\textbf{Par de tipos} & \textbf{SD 1.5} & \textbf{SD 2.1} & \textbf{SDXL} \\
\midrule
Color vs Iluminación    & $0{,}093 \pm 0{,}048$ & $0{,}083 \pm 0{,}052$ & $0{,}119 \pm 0{,}086$ \\
Color vs Estilo         & $0{,}056 \pm 0{,}050$ & $0{,}051 \pm 0{,}039$ & $0{,}227 \pm 0{,}121$ \\
Iluminación vs Estilo   & $0{,}068 \pm 0{,}058$ & $0{,}087 \pm 0{,}079$ & $0{,}328 \pm 0{,}130$ \\
\bottomrule
\end{tabular}
\end{table}

Existe una \textbf{tensión entre consistencia y disentanglement}: SD~1.5 y SD~2.1 presentan excelente ortogonalidad ($|\cos| < 0{,}10$) pero sus direcciones son dependientes del contexto; SDXL invierte este equilibrio con mayor consistencia pero mayor entrelazamiento, especialmente entre iluminación y estilo ($|\cos| = 0{,}328$).

\subsection{Alineación de las direcciones semánticas con la base PCA}

\begin{table}[htb]
\centering
\caption{Fracción media de energía de las direcciones semánticas en los diez primeros componentes principales.}
\label{tab:energia_pca}
\small
\begin{tabular}{lrrr}
\toprule
\textbf{Tipo} & \textbf{SD 1.5} & \textbf{SD 2.1} & \textbf{SDXL} \\
\midrule
Color       & 0,72 & 0,67 & \textbf{0,95} \\
Iluminación & 0,56 & 0,50 & \textbf{0,93} \\
Estilo      & 0,51 & 0,51 & \textbf{0,79} \\
\bottomrule
\end{tabular}
\end{table}

En SDXL, el 95\% de la energía de las direcciones de color se concentra en los diez primeros componentes, con el PC0 sistemáticamente alineado con las direcciones de color y el PC3 con las de iluminación, permitiendo interpretar estas componentes como \textit{dimensiones de color} y \textit{dimensión de iluminación} del espacio latente de SDXL.

\section{Fine-tuning y deriva del espacio latente}
\label{sec:res:finetune}

\subsection{Deriva geométrica tras el fine-tuning}

La Tabla~\ref{tab:deriva_geometrica} resume las métricas de deriva geométrica obtenidas regenerando el corpus completo con cada modelo adaptado y comparando los latentes resultantes con los del pretrain para los mismos prompts y semillas. Las Figuras~\ref{fig:deriva_sd15}, \ref{fig:deriva_sd21} y \ref{fig:deriva_sdxl} muestran el análisis detallado por modelo en seis paneles: distribuciones de distancias L2 intra/inter-prompt antes y después del fine-tuning (izquierda), separabilidad y participation ratio normalizados (centro) y distribución del desplazamiento por latente (derecha) para la geometría (fila superior); consistencia semántica intra-tipo, energía PCA de las direcciones de atributo y curva de alineación de Grassmann de la base PCA (fila inferior).

\begin{table}[htb]
\centering
\caption{Métricas de deriva geométrica pre/post fine-tuning WikiArt. $\bar{\delta} = \|\z_0^\text{ft} - \z_0^\text{pre}\|_2$ medio por latente; $\cos_c$ es la similitud coseno entre centroides; KS es el estadístico del test Kolmogorov--Smirnov sobre las distribuciones de distancias inter-prompt (valor $p$ entre paréntesis).}
\label{tab:deriva_geometrica}
\small
\begin{tabular}{lrrrrrrrr}
\toprule
\textbf{Modelo} & \textbf{sep\textsubscript{pre}} & \textbf{sep\textsubscript{ft}} & $\Delta\text{sep}$ & \textbf{PR\textsubscript{pre}} & \textbf{PR\textsubscript{ft}} & $\Delta\text{PR}$ & $\cos_c$ & $\bar{\delta}$ \\
\midrule
SD 1.5 & 0,976 & 0,974 & $-0{,}001$ & 12,35 & 11,68 & $-0{,}67$ & 0,991 & \phantom{0}36,6 \\
SD 2.1 & 0,966 & 0,959 & $-0{,}007$ & 12,33 & 10,75 & $-1{,}58$ & 0,993 & \phantom{0}39,6 \\
SDXL   & 1,118 & 1,013 & $\mathbf{-0{,}105}$ & 11,56 & 10,11 & $-1{,}45$ & 0,914 & $\mathbf{138,4}$ \\
\bottomrule
\end{tabular}
\end{table}

SD~1.5 y SD~2.1 muestran una deriva geométrica prácticamente nula. El ratio de separabilidad varía en menos de 0,01 puntos y los desplazamientos medios por latente ($\bar{\delta} = 36{,}6$ y $39{,}6$) representan apenas el 24\% y el 20\% de las distancias intra-prompt pretrain. Los centroides se mantienen casi idénticos ($\cos_c > 0{,}99$) y el test KS no rechaza la igualdad de distribuciones de distancias inter-prompt ($p = 0{,}182$ y $p = 0{,}665$, respectivamente).

SDXL presenta un perfil cualitativamente distinto. El desplazamiento medio por latente es de 138,4, equivalente al 53\% de la distancia intra-prompt pretrain ($\mu_\text{intra} \approx 261$): los latentes regenerados de SDXL se han desplazado más de la mitad de la distancia típica entre semillas del mismo prompt. El ratio de separabilidad cae 0,105 puntos, perdiendo casi toda la ventaja de separabilidad característica de SDXL pretrain (sep = 1,118 $\to$ 1,013). La similitud coseno entre centroides cae a 0,914, registrando la mayor desviación direccional de los tres modelos. A pesar de esta magnitud, el test KS tampoco rechaza la igualdad de distribuciones ($p = 0{,}120$), indicando que la forma global de la distribución de distancias se preserva aunque su posición varíe.

\begin{figure}[htbp]
\centering
\includegraphics[width=\textwidth]{data/fase5/deriva_sd15}
\caption{Deriva latente pre/post fine-tuning WikiArt para SD~1.5. La distribución de distancias, la separabilidad y el PR apenas varían; el desplazamiento medio por latente es de 36,6. La consistencia semántica y la energía PCA de las direcciones se mantienen estables. La base PCA conserva una alineación de Grassmann de 0,82 en los diez primeros componentes.}
\label{fig:deriva_sd15}
\end{figure}

\begin{figure}[htbp]
\centering
\includegraphics[width=\textwidth]{data/fase5/deriva_sd21}
\caption{Deriva latente pre/post fine-tuning WikiArt para SD~2.1. La geometría es prácticamente estable. La consistencia de estilo muestra el único incremento apreciable entre todos los modelos ($0{,}010 \to 0{,}023$, +130\%). La alineación PCA top-10 es la más alta de los tres modelos (0,888).}
\label{fig:deriva_sd21}
\end{figure}

\begin{figure}[htbp]
\centering
\includegraphics[width=\textwidth]{data/fase5/deriva_sdxl}
\caption{Deriva latente pre/post fine-tuning WikiArt para SDXL. El eje de desplazamiento por latente alcanza 250 frente a 100--110 en SD~1.5 y SD~2.1. La separabilidad cae desde 1,118 hasta 1,013. La consistencia semántica disminuye en todos los atributos (color: $-22\%$, iluminación: $-52\%$, estilo: $-37\%$). La alineación PCA top-10 cae a 0,537.}
\label{fig:deriva_sdxl}
\end{figure}

\subsection{Rotación de la base PCA tras el fine-tuning}

\begin{table}[htb]
\centering
\caption{Alineación de Grassmann entre los subespacios PCA pretrain y post fine-tuning. El valor 1 indica subespacios idénticos; valores menores indican rotación de la base.}
\label{tab:pca_alineacion}
\small
\begin{tabular}{lrr}
\toprule
\textbf{Modelo} & \textbf{Alineación top-10} & \textbf{Alineación top-50} \\
\midrule
SD 1.5 & 0,817 & 0,725 \\
SD 2.1 & \textbf{0,888} & \textbf{0,800} \\
SDXL   & 0,537 & 0,282 \\
\bottomrule
\end{tabular}
\end{table}

La alineación de Grassmann cuantifica cuánto ha rotado la base de varianza principal (Tabla~\ref{tab:pca_alineacion}). SD~1.5 y SD~2.1 mantienen alineaciones de 0,82--0,89 en los diez primeros componentes, indicando que el fine-tuning no reorganiza significativamente las direcciones de mayor varianza. En SDXL, la alineación cae a 0,537 en los diez primeros componentes y a 0,282 en los cincuenta primeros: el subespacio de máxima varianza post fine-tuning comparte apenas la mitad de la orientación del subespacio pretrain para los componentes principales y menos de un tercio para el subespacio extendido.

\subsection{Deriva semántica de las direcciones de atributo}

\begin{table}[htb]
\centering
\caption{Consistencia intra-tipo pre y post fine-tuning WikiArt. Los valores entre paréntesis indican el cambio porcentual relativo.}
\label{tab:deriva_semantica_consistencia}
\small
\begin{tabular}{lcccccc}
\toprule
\textbf{Tipo} & \multicolumn{2}{c}{\textbf{SD 1.5}} & \multicolumn{2}{c}{\textbf{SD 2.1}} & \multicolumn{2}{c}{\textbf{SDXL}} \\
\cmidrule(lr){2-3}\cmidrule(lr){4-5}\cmidrule(lr){6-7}
 & Pre & Post & Pre & Post & Pre & Post \\
\midrule
Color       & 0,300 & 0,276 {\small($-8\%$)}  & 0,229 & 0,231 {\small($+1\%$)}   & 0,707 & 0,554 {\small($-22\%$)} \\
Iluminación & 0,109 & 0,116 {\small($+6\%$)}  & 0,155 & 0,163 {\small($+5\%$)}   & 0,517 & 0,247 {\small($-52\%$)} \\
Estilo      & 0,028 & 0,027 {\small($-4\%$)}  & 0,010 & 0,023 {\small($+130\%$)} & 0,209 & 0,132 {\small($-37\%$)} \\
\bottomrule
\end{tabular}
\end{table}

La Tabla~\ref{tab:deriva_semantica_consistencia} recoge el resultado más llamativo de la Fase~5. SD~1.5 y SD~2.1 exhiben estabilidad semántica casi total: los cambios en la consistencia intra-tipo son inferiores al 8\% relativo en todos los atributos, con la excepción de la consistencia de estilo en SD~2.1, que pasa de 0,010 a 0,023 (+130\% relativo). Este incremento, aunque grande en términos porcentuales, parte de un valor de base prácticamente nulo y el resultado post fine-tuning (0,023) sigue siendo inferior en un orden de magnitud a la consistencia de estilo pretrain de SDXL (0,209).

SDXL muestra una deriva semántica severa en todos los atributos. La consistencia de color cae un 22\% (0,707 $\to$ 0,554), la de iluminación un 52\% (0,517 $\to$ 0,247) y la de estilo un 37\% (0,209 $\to$ 0,132). Los valores post fine-tuning de SDXL siguen siendo superiores a los pretrain de SD~1.5 y SD~2.1, confirmando que la representación semántica de SDXL, aunque perturbada, mantiene ventaja cualitativa. La disminución más severa corresponde a la iluminación, atributo cuya dirección estaba parcialmente alineada con el estilo antes del fine-tuning (Tabla~\ref{tab:ortogonalidad}), sugiriendo que la reorganización inducida por WikiArt afecta especialmente a la representación de propiedades con correlación visual en el dominio artístico.

\subsection{Comparativa cross-modelo de la deriva latente}

\begin{figure}[htbp]
\centering
\includegraphics[width=\textwidth]{data/deriva_cross_model}
\caption{Comparativa cross-modelo de la deriva latente tras fine-tuning WikiArt. Fila superior (izq.\ a der.): $\Delta$sep, $\Delta$PR y alineación PCA top-10. Fila inferior: $\Delta$consistencia intra-tipo, $\Delta$energía PCA de las direcciones semánticas y desplazamiento latente medio $\bar{\delta}$. El modelo SDXL (verde) domina todas las métricas de deriva geométrica y semántica.}
\label{fig:deriva_cross}
\end{figure}

La Figura~\ref{fig:deriva_cross} sintetiza la comparativa entre los tres modelos. El patrón es uniforme y consistente: SD~1.5 y SD~2.1 presentan deltas cercanos a cero en prácticamente todas las métricas, mientras que SDXL exhibe las mayores magnitudes de cambio en separabilidad ($\Delta\text{sep} = -0{,}105$, frente a $-0{,}001$ y $-0{,}007$), desplazamiento latente medio ($\bar{\delta} = 138{,}4$, frente a $36{,}6$ y $39{,}6$), alineación PCA top-10 (0,537, frente a 0,817 y 0,888), y consistencia semántica de color ($-0{,}153$) e iluminación ($-0{,}270$). La notable excepción es el $\Delta\text{PR}$, donde SD~2.1 registra la mayor pérdida de dimensionalidad efectiva ($-1{,}58$), ligeramente superior a la de SDXL ($-1{,}45$), mientras que SD~1.5 pierde solo $-0{,}67$ dimensiones efectivas. Este patrón sugiere que SD~2.1, aunque geométricamente estable a nivel de distancias, experimenta la mayor concentración dimensional de varianza tras el fine-tuning.

% ===========================================================
%  DISCUSIÓN
% ===========================================================
\chapter{Discusión}
\label{ch:discusion}

\section{Convergencia geométrica y divergencia semántica: el patrón fundamental}
\label{sec:dis:patron}

El hallazgo central de este trabajo es la coexistencia de dos tendencias opuestas según el nivel de análisis. A nivel geométrico global, los tres modelos convergen: el participation ratio PR $\approx 12$ es prácticamente idéntico pese a que las dimensiones nominales difieren en un factor de cuatro (Tabla~\ref{tab:dimensionalidad}). A nivel semántico, los modelos divergen de forma cualitativa: la consistencia de las direcciones de atributo en SDXL supera entre 2,3 y 7,5 veces los valores de SD~1.5 y SD~2.1 (Tabla~\ref{tab:consistencia}). Esta asimetría no era predecible a priori y constituye la aportación conceptual más relevante del trabajo.

La interpretación más parsimoniosa es que la baja dimensionalidad intrínseca es una propiedad del \textit{proceso de difusión}, no de la arquitectura específica. El proceso de denoising, con independencia de la capacidad del UNet y del codificador de texto, colapsa la distribución de latentes $\z_0$ sobre una variedad de dimensión intrínseca aproximadamente constante. En cambio, la organización semántica de esa variedad ---la manera en que los atributos visuales quedan representados dentro de ella--- depende críticamente de la calidad del condicionamiento textual.

Esta disociación tiene una consecuencia metodológica directa: los instrumentos que operan sobre la geometría global del espacio $\z_0$ (PCA, participation ratio, correlación con métricas de imagen) son transferibles entre arquitecturas y ofrecen resultados comparables. Los instrumentos que operan sobre la organización semántica (consistencia de direcciones, alineación PCA-semántica) no lo son: sus resultados dependen del codificador de texto y no se pueden extrapolar de una arquitectura a otra.

\section{La dimensionalidad intrínseca como propiedad del proceso de difusión}
\label{sec:dis:dimensionalidad}

Los valores PR $\approx 12$ son significativamente menores que las 40--60 dimensiones efectivas reportadas por Pope et al.\ (2021) para el espacio $\mathcal{W}$ de StyleGAN con un corpus de rostros humanos. Esta diferencia admite dos explicaciones complementarias.

La primera es el tipo de corpus analizado. Los 60 latentes base cubren 20 dominios visuales distintos con solo 3 muestras por dominio. La dimensionalidad efectiva del PCA refleja en este régimen cuántos ejes independientes de variación semántica existen entre los 20 dominios, no la riqueza interna de cada uno. Doce dimensiones para describir la diversidad de 20 escenas heterogéneas es un valor plausible y consistente con la estructura del corpus. Pope et al.\ trabajaron con miles de rostros cuya diversidad intra-dominio es mucho mayor.

La segunda explicación es de naturaleza arquitectónica. La red de mapeo de StyleGAN aprende a transformar el espacio esférico $\mathcal{Z}$ en un espacio $\mathcal{W}$ con mayor disentanglement, un proceso que puede expandir la dimensionalidad efectiva. El proceso de denoising de los LDMs actúa como un proyector implícito sobre la variedad de imágenes naturales; la baja dimensionalidad del espacio $\z_0$ puede ser consecuencia de que este proyector es más restrictivo.

Lo más informativo no es el valor absoluto de PR $\approx 12$ sino su invarianza entre modelos. El UNet de SDXL tiene tres veces más parámetros que el de SD~1.5 y su dimensión nominal latente es cuatro veces mayor, pero el PR no varía. Esto implica que ni la capacidad del UNet ni el tamaño del espacio latente determinan cuántas dimensiones efectivas ocupa la distribución de $\z_0$: la dimensionalidad intrínseca está fijada por la distribución de prompts del corpus y por la naturaleza del proceso de difusión, no por la arquitectura. Esta predicción es verificable: un corpus con mayor diversidad semántica debería producir un PR más alto, igualmente invariante entre modelos.

\section{El codificador de texto como modulador semántico del espacio latente}
\label{sec:dis:text_encoder}

La comparación entre SD~1.5 y SD~2.1 actúa como control experimental parcial: ambos modelos emplean un único codificador de texto pero de distinta capacidad (CLIP ViT-L/14 con 768 dimensiones frente a OpenCLIP ViT-H/14 con 1.024 dimensiones). Los resultados muestran que este incremento del 33\% en la dimensión del embedding no produce ningún cambio cualitativo en la consistencia intra-tipo (Tabla~\ref{tab:consistencia}): los valores permanecen en el rango 0,01--0,30 para ambos modelos sin diferencias sistemáticas. La dirección de estilo en SD~2.1 tiene incluso menor consistencia (0,010) que en SD~1.5 (0,028), descartando una tendencia monotónica.

Este resultado tiene una implicación fundamental: \textit{escalar un único codificador CLIP no es suficiente para inducir representaciones semánticas universales en el espacio $\z_0$}. El salto cualitativo se produce únicamente con SDXL, que introduce un segundo codificador de la familia OpenCLIP (OpenCLIP ViT-bigG/14, 1.280 millones de parámetros) en modo de condicionamiento global, adicionalmente al CLIP ViT-L/14 de SD~1.5. La consistencia de color salta de 0,229 a 0,707 (factor $\times 3$), la de iluminación de 0,155 a 0,517 (factor $\times 3{,}3$) y la de estilo de 0,010 a 0,209 (factor $\times 21$).

Una hipótesis explicativa compatible con estos datos es que el condicionamiento dual de SDXL introduce una forma de anclaje semántico que los modelos de codificador único no poseen. CLIP ViT-L/14 aporta información semántica fina mediante atención cruzada localizada en el UNet, mientras que OpenCLIP ViT-bigG/14 aporta un embedding global de la escena completa que condiciona el comportamiento del UNet de forma más holística. La co-presencia de estas dos señales de condicionamiento con alcances semánticos distintos podría guiar al UNet a aprender representaciones $\z_0$ donde los atributos visuales se alinean de forma más consistente con las instrucciones textuales. Esta redundancia informativa ---dos descripciones del mismo prompt a distintas escalas--- sería el mecanismo que ancla la representación semántica.

Esta hipótesis no puede confirmarse de forma definitiva con los datos disponibles, ya que los tres modelos difieren simultáneamente en codificador de texto, capacidad del UNet, resolución nativa, datos de entrenamiento y cómputo de preentrenamiento. Sin embargo, el hecho de que SD~2.1 ---con un codificador individual de mayor capacidad que el de SD~1.5 y un UNet de capacidad similar--- no muestre mejora cualitativa hace del codificador dual la variable diferencial más parsimoniosa. La limitación causal de este argumento se discute en la sección \ref{sec:dis:causalidad}.

\section{La tensión entre consistencia y disentanglement}
\label{sec:dis:tension}

El patrón de ortogonalidad inter-tipo revela una tensión que no ha sido documentada en la literatura previa: las arquitecturas que producen direcciones semánticas más universales (alta consistencia intra-tipo) presentan también mayor entrelazamiento entre tipos (Tabla~\ref{tab:ortogonalidad}). En SDXL, el par iluminación-estilo alcanza una similitud coseno absoluta media de 0,328, frente a valores inferiores a 0,10 en SD~1.5 y SD~2.1.

Esta tensión puede interpretarse desde dos perspectivas no excluyentes. Desde una perspectiva de capacidad representacional, SD~1.5 y SD~2.1 tienen un presupuesto semántico insuficiente para mantener direcciones universales: sus vectores de atributo son locales y específicos de la escena porque el condicionamiento textual no es suficientemente potente para imponer una dirección consistente a través de distintos dominios visuales. La ortogonalidad inter-tipo es parcialmente trivial en este caso: si las direcciones son esencialmente ruido orientado por la geometría local de la escena, la correlación entre pares de tipos será naturalmente baja.

Desde una perspectiva de correlación semántica, el entrelazamiento de iluminación y estilo en SDXL puede reflejar correlaciones genuinas del mundo visual. Los estilos artísticos poseen tratamientos lumínicos característicos: el claroscuro del barroco, la luz difusa del impresionismo, el contraste saturado del expresionismo. Un modelo con suficiente capacidad semántica capturará estas correlaciones intrínsecas, y el entrelazamiento sería entonces no un defecto sino una representación más fiel del espacio visual. Bajo esta interpretación, intentar eliminar el efecto sobre el estilo al editar la iluminación sería antinatural y produciría artefactos.

Para el diseño de sistemas de edición, esta distinción importa operativamente. Si el entrelazamiento es un artefacto del entrenamiento, puede mitigarse proyectando la dirección de edición sobre el subespacio ortogonal al complementario. Si es semánticamente informativo, esa proyección empeoraría el resultado visual. Discriminar entre estas dos hipótesis requeriría estudios de evaluación perceptual humana que están fuera del alcance de este trabajo.

\section{La separabilidad como índice de confianza semántica}
\label{sec:dis:separabilidad}

El ratio de separabilidad sep $< 1$ en SD~1.5 (0,976) y SD~2.1 (0,966) tiene una interpretación directa con consecuencias prácticas: en estos modelos, la variabilidad estocástica del proceso de denoising produce desplazamientos en el espacio $\z_0$ de magnitud comparable o mayor a los producidos por un cambio de contenido semántico entre escenas distintas. Regenerar la misma escena con otra semilla desplaza el latente tanto como pasar a generar una escena completamente diferente.

Esta propiedad implica una relación señal-ruido desfavorable para las técnicas de edición basadas en aritmética de vectores en $\z_0$ con SD~1.5 y SD~2.1. Para extraer una dirección semántica limpia es necesario promediar sobre múltiples semillas, procedimiento que el diseño del corpus ya incorpora con las tres semillas por prompt. Aun así, la magnitud de la señal semántica neta es modesta, y las ediciones basadas en $\z_0$ en estos modelos tienen una variabilidad residual alta que puede traducirse en cambios no deseados.

SDXL invierte esta relación (sep $= 1{,}118$, $p < 0{,}001$): el contenido del prompt es sistemáticamente el predictor dominante de la posición del latente, y la variabilidad estocástica queda subordinada al contenido semántico. Esta propiedad explica, al menos parcialmente, la mayor consistencia intra-tipo de las direcciones semánticas: si la posición del latente está determinada principalmente por el contenido del prompt, la diferencia entre dos versiones del mismo prompt con atributo modificado produce un desplazamiento más coherente y predecible.

La conexión con el trabajo de Kwon et al.\ (2023) es directa: si el h-space ---las activaciones del bottleneck del UNet--- exhibe mayor linealidad semántica en SDXL que en SD~1.5, es coherente esperar que el espacio $\z_0$ resultante del proceso completo de denoising también exhiba mayor organización semántica. La separabilidad y la consistencia son manifestaciones del mismo fenómeno medidas, respectivamente, como organización del espacio de salida y como coherencia de las representaciones de atributo.

\section{Hacia un GANSpace aplicado a modelos de difusión latente}
\label{sec:dis:ganspace}

El hallazgo de que en SDXL el 95\% de la energía de las direcciones de color y el 93\% de la de iluminación se concentran en los diez primeros componentes principales (Tabla~\ref{tab:energia_pca}), con PC0 alineado con el color y PC3 con la iluminación, tiene una implicación metodológica de primer orden: para SDXL, el PCA sobre un conjunto representativo de latentes $\z_0$ es una forma de descubrimiento no supervisado de direcciones semánticas de atributo.

Esta propiedad establece un puente directo con la metodología GANSpace (Härkönen et al., 2020), que demostró que el PCA sobre vectores $\mathbf{w}$ de StyleGAN produce componentes semánticamente interpretables sin necesitar etiquetas ni supervisión. El presente resultado extiende esa propiedad al espacio $\z_0$ de SDXL, con la ventaja añadida de que $\z_0$ es directamente accesible mediante inversión DDIM sin recorrer ninguna red de mapeo explícita. Un usuario que disponga de una colección de imágenes y desee construir un vocabulario semántico básico puede obtenerlo aplicando PCA a los $\z_0$ invertidos de SDXL, encontrando que los primeros componentes tienen interpretabilidad visual directa.

En SD~1.5 y SD~2.1, la fracción de energía semántica en los diez primeros componentes es del 50--72\%, insuficiente para una alineación clara y consistente. Los componentes PCA reflejan en estos modelos principalmente la varianza estructural de las escenas, no la organización de atributos visuales. Para SD~1.5 y SD~2.1, el enfoque supervisado de InterFaceGAN (Shen et al., 2020) ---encontrar la dirección mediante un clasificador lineal sobre latentes etiquetados--- es el método más adecuado. La diferencia entre modelos en este aspecto no es solo cuantitativa sino cualitativa: el umbral a partir del cual el PCA descubre semántica parece situarse en algún punto entre SD~2.1 y SDXL, coincidiendo con el salto en el condicionamiento textual.

\section{El fine-tuning en WikiArt: resultados frente a hipótesis}
\label{sec:dis:finetune}

Los resultados del análisis pretrain generaron tres predicciones específicas, formuladas al final del análisis de las propiedades semánticas de cada arquitectura. Contrastarlas con los resultados experimentales de la Fase~5 (sección~\ref{sec:res:finetune}) es el principal hilo conductor de esta sección.

\textbf{Predicción 1 refutada: SDXL no amplifica su consistencia semántica; la pierde.} La hipótesis era que un modelo con representación de estilo ya parcialmente universal absorbería el conocimiento artístico de forma más eficiente, amplificando la consistencia. El resultado es el opuesto: SDXL pierde el 37\% de su consistencia de estilo (0,209 $\to$ 0,132), el 52\% de la de iluminación (0,517 $\to$ 0,247) y el 22\% de la de color (0,707 $\to$ 0,554). SD~1.5 y SD~2.1 permanecen prácticamente inalterados, con la única excepción de SD~2.1, que incrementa su consistencia de estilo de 0,010 a 0,023 (+130\% relativo, pero en valor absoluto insignificante). La dirección del efecto es la inversa a la predicha.

\textbf{Predicción 2 refutada: SDXL sufre la mayor deriva geométrica.} La hipótesis era que SD~1.5 y SD~2.1 requerirían cambios más profundos en su geometría al carecer de representación coherente de estilo. Los datos muestran exactamente lo contrario: el desplazamiento latente medio de SDXL ($\bar{\delta} = 138{,}4$) es 3,5 veces superior al de SD~1.5 y SD~2.1 ($\bar{\delta} \approx 37$--$40$), la similitud coseno entre centroides cae a 0,914 frente a 0,991--0,993, y la alineación de la base PCA se reduce a 0,537 frente a 0,817--0,888 para los otros dos modelos.

\textbf{Interpretación: la fragilidad semántica de las representaciones organizadas.} La coherencia de estos dos resultados sugiere un principio que no está documentado en la literatura de personalización: los modelos con espacios latentes más organizados semánticamente en pretrain son más susceptibles a la disrupción por adaptación de dominio, no menos. Una posible explicación es que la representación semántica estructurada de SDXL está fuertemente anclada en el régimen visual del preentrenamiento (fotografía naturalista de LAION-5B). Cuando el fine-tuning impone un dominio cualitativamente distinto ---pintura clásica, con una distribución de textura, paleta y estructura espacial opuesta a la fotografía--- el UNet de SDXL debe reestructurar sus representaciones internas de forma más drástica que los modelos cuyos espacios $\z_0$ eran ya localmente caóticos. La plasticity de SD~1.5 y SD~2.1 ante el fine-tuning es consecuencia directa de su menor organización pretrain.

Esta interpretación tiene una implicación práctica directa: para aplicaciones de edición basadas en aritmética de vectores en $\z_0$, el fine-tuning sobre WikiArt degrada significativamente la fiabilidad de SDXL como plataforma, mientras que SD~1.5 y SD~2.1 mantienen sus propiedades (modestas) intactas. SDXL tras fine-tuning WikiArt tiene las ventajas de ninguno de los dos estados: menor consistencia semántica que el pretrain de SDXL y mayor que el pretrain de SD~1.5/SD~2.1, con una geometría reorientada que invalida las direcciones de atributo calculadas sobre el modelo base.

\textbf{La congelación del codificador de texto como garantía de interpretabilidad.} La congelación del codificador de texto durante el fine-tuning es el elemento de diseño experimental que hace interpretables estos resultados. Garantiza que cualquier cambio en el espacio $\z_0$ sea atribuible exclusivamente a la reorganización del proceso de denoising del UNet, no a cambios en el espacio de condicionamiento textual. Esta separación experimental no ha sido abordada sistemáticamente en la literatura de personalización (Gal et al., 2022; Ruiz et al., 2023), donde el UNet y el codificador de texto se modifican frecuentemente de forma simultánea o se ajusta solo el codificador, sin aislar sus contribuciones respectivas. Bajo este diseño, los resultados confirman que la adaptación de dominio actúa sobre la dinámica de denoising del UNet, modificando el punto del espacio latente al que converge el proceso, no el significado semántico de los tokens de condicionamiento.

\section{Posicionamiento en la literatura y brechas cubiertas}
\label{sec:dis:literatura}

Respecto a Kwon et al.\ (2023), este trabajo complementa su análisis del h-space identificando propiedades análogas en el espacio $\z_0$: la linealidad semántica en $\z_0$ es más marcada para SDXL, lo que sugiere que la arquitectura del UNet modula no solo las activaciones intermedias sino también el espacio latente final. Crucialmente, este trabajo cubre las tres principales arquitecturas de la familia SD, mientras que Kwon et al.\ se limitan a SD~1.x sin comparación sistemática.

Respecto a Pope et al.\ (2021), los valores PR $\approx 12$ son menores que las 40--60 dimensiones reportadas para el $\mathcal{W}$ de StyleGAN para rostros, por razones discutidas en la sección \ref{sec:dis:dimensionalidad}. La invarianza del PR entre arquitecturas es un resultado nuevo: Pope et al.\ estudian un único modelo y no tienen hipótesis sobre transferibilidad arquitectónica.

Respecto a los trabajos de personalización (Gal et al., 2022; Ruiz et al., 2023; Hu et al., 2022), que evalúan el fine-tuning en términos de calidad visual sin caracterizar la estructura del espacio latente resultante, este trabajo proporciona la línea de base cuantitativa que permite anticipar el impacto del fine-tuning a nivel estructural, cubriendo la tercera brecha identificada en la revisión del estado del arte.

La cuarta brecha---la ausencia de un protocolo replicable para la evaluación comparativa del espacio latente de LDMs aplicable uniformemente a distintas arquitecturas---queda cubierta por el protocolo desarrollado en este trabajo, cuya transferibilidad a cualquier LDM con VAE de factor de compresión $f=8$ está garantizada por diseño.

\section{Alcance causal de la interpretación arquitectónica}
\label{sec:dis:causalidad}

Una limitación epistemológica de fondo es que la comparación entre modelos es observacional, no experimental en sentido estricto. Los tres modelos difieren simultáneamente en el codificador de texto, la capacidad del UNet, la resolución nativa, los datos de entrenamiento y el cómputo total de preentrenamiento. La atribución causal de las diferencias observadas al codificador de texto dual ---aunque es la hipótesis más parsimoniosa, argumentada en la sección \ref{sec:dis:text_encoder}--- no puede establecerse de forma definitiva.

El experimento de control ideal sería mantener fijo el UNet de SDXL y variar únicamente el codificador de texto: entrenar el mismo UNet con el codificador doble, con el codificador de SD~1.5 y con el de SD~2.1, sobre los mismos datos y con el mismo cómputo. Un experimento de ese tipo no existe en la literatura pública y excede las capacidades computacionales de este trabajo.

En ausencia de ese control, los resultados son más útiles como caracterización empírica de las arquitecturas disponibles que como validación de una teoría causal. La regularidad observada ---PR invariante entre arquitecturas, consistencia dependiente del codificador de texto, tensión consistencia-disentanglement--- es suficientemente robusta para guiar decisiones de diseño práctico (qué arquitectura usar para qué tarea de edición) y para generar las hipótesis verificables sobre los resultados del fine-tuning formuladas en la sección \ref{sec:dis:finetune}.

% ===========================================================
%  CONCLUSIONES
% ===========================================================
\chapter{Conclusiones}
\label{ch:conclusiones}

\section{Síntesis de las contribuciones}
\label{sec:concl:contribuciones}

Este trabajo ha presentado el primer análisis comparativo sistemático y cuantitativo de la estructura del espacio latente $\z_0$ de tres arquitecturas de la familia Stable Diffusion bajo control de la variable VAE. Las contribuciones principales son cuatro.

\textbf{Primera contribución: protocolo de evaluación estandarizado.} Se ha desarrollado y validado un protocolo reproducible para caracterizar el espacio latente de cualquier LDM que comparta la arquitectura VAE de Stable Diffusion. El protocolo abarca la generación del corpus de latentes, el análisis geométrico (separabilidad, dimensionalidad efectiva, coherencia perceptual) y el análisis semántico (discriminabilidad, consistencia, ortogonalidad y alineación PCA de las direcciones de atributo). La homogeneidad de prompts, semillas e hiperparámetros entre modelos garantiza que las diferencias observadas son comparables sin artefactos metodológicos.

\textbf{Segunda contribución: baja dimensionalidad intrínseca universal.} El participation ratio PR $\approx 12$ es prácticamente idéntico en los tres modelos con independencia de sus dimensiones nominales (16.384, 36.864 y 65.536 elementos). Esta convergencia, discutida en la sección \ref{sec:dis:dimensionalidad}, justifica la aplicación de técnicas lineales como el PCA y las interpolaciones en el espacio $\z_0$ incluso para arquitecturas de muy distinto tamaño, y señala al proceso de difusión ---no a la arquitectura--- como la fuente de la baja dimensionalidad.

\textbf{Tercera contribución: el codificador de texto como factor determinante de la universalidad semántica.} La transición al doble codificador de SDXL produce un salto cualitativo en la consistencia de las direcciones semánticas que no se obtiene escalando un único codificador CLIP: la consistencia del color pasa de 0,30 a 0,71, la iluminación de 0,11 a 0,52 y el estilo de 0,03 a 0,21. Este resultado, interpretado en la sección \ref{sec:dis:text_encoder}, no puede explicarse únicamente por el mayor tamaño del UNet.

\textbf{Cuarta contribución: identificación de la tensión entre consistencia y disentanglement.} Los modelos más pequeños presentan excelente ortogonalidad inter-tipo ($|\cos| < 0{,}10$) pero sus direcciones semánticas son contexto-dependientes. SDXL invierte este equilibrio con mayor consistencia pero mayor entrelazamiento. Este patrón, analizado en la sección \ref{sec:dis:tension}, tiene consecuencias directas para el diseño de sistemas de edición controlada.

\textbf{Quinta contribución: fragilidad semántica bajo adaptación de dominio.} El fine-tuning LoRA sobre WikiArt revela un patrón inverso al esperado: SDXL, con la representación pretrain más organizada, sufre la mayor perturbación geométrica y semántica ($\bar{\delta} = 138{,}4$, alineación PCA = 0,54, pérdidas de consistencia del 22--52\%), mientras que SD~1.5 y SD~2.1 muestran deriva prácticamente nula. Este resultado, discutido en la sección \ref{sec:dis:finetune}, establece que una mayor organización semántica pretrain no protege frente a la disrupción por adaptación de dominio, sino que la amplifica. Tiene implicaciones directas para la selección de arquitectura en aplicaciones de personalización donde se requiera preservar las propiedades de edición del espacio latente.

\section{Limitaciones}
\label{sec:concl:limitaciones}

El trabajo presenta cuatro limitaciones principales. \textbf{Tamaño del corpus}: 60 latentes base y 90 de atributo por modelo es suficiente para la metodología implementada pero limita la aplicación de estimadores no paramétricos de dimensionalidad intrínseca como TwoNN y la resolución del análisis PCA. \textbf{Tres atributos visuales}: el análisis no cubre atributos geométricos, relacionales ni temporales, que pueden mostrar patrones distintos a los aquí documentados. \textbf{Restricciones computacionales}: la GPU T4 de Colaboratory impone límites en el número de pasos de fine-tuning de SDXL (800 pasos frente a 1.000 de los modelos menores), lo que podría subestimar la deriva latente definitiva de SDXL o, por el contrario, no alcanzar el punto de mínimo de pérdida. \textbf{Un único dominio de fine-tuning}: los resultados de la Fase~5 corresponden exclusivamente al dominio de pintura clásica (WikiArt). No es posible determinar si la fragilidad semántica observada en SDXL es específica de este dominio artístico ---cualitativamente opuesto a la fotografía naturalista de LAION-5B--- o generalizable a otros dominios de adaptación.

\section{Trabajo futuro}
\label{sec:concl:futuro}

\textbf{Extensión a arquitecturas DiT.} Los modelos SD3 y Flux utilizan transformadores de difusión. Verificar si la baja dimensionalidad intrínseca y las propiedades de separabilidad semántica se mantienen en estas arquitecturas es directamente relevante para evaluar la transferibilidad de los resultados al estado del arte actual.

\textbf{Comparación de $\z_0$ con el h-space.} Un análisis sistemático que aplique las mismas métricas de consistencia, ortogonalidad y alineación PCA tanto al espacio $\z_0$ como a las activaciones del h-space para los tres modelos permitiría determinar cuál de los dos espacios es más apropiado para distintas tareas de edición.

\textbf{Corpus extendido y estimadores no paramétricos.} Un corpus de 500--1.000 muestras por modelo permitiría aplicar el estimador TwoNN para una medida más robusta de la dimensionalidad intrínseca y, según la hipótesis de la sección \ref{sec:dis:dimensionalidad}, verificar si el PR crece con la diversidad semántica del corpus manteniéndose invariante entre modelos.

\textbf{Métricas perceptuales profundas.} La sustitución del MSE y el SSIM global por LPIPS en el análisis de correlación proporcionaría una medida más relevante perceptualmente y, coherentemente con el argumento sobre el rol del VAE, una correlación más uniforme entre modelos.

\textbf{Fine-tuning con distintos dominios.} Replicar el experimento con dominios de naturaleza distinta ---fotografía médica, imágenes satelitales, ilustración de producto--- permitiría determinar si las diferencias de respuesta entre arquitecturas son específicas del dominio artístico o generalizables.

\textbf{Análisis espectral de los pesos LoRA.} El análisis de los valores singulares de las matrices $\mathbf{B}\mathbf{A}$ por capa permitiría identificar qué partes del UNet concentran el mayor cambio durante el fine-tuning y correlacionar esa localización con los cambios en el espacio $\z_0$.

\section{Reflexión final}
\label{sec:concl:reflexion}

Este trabajo parte de una pregunta deceptivamente simple: ¿son los espacios latentes de SD~1.5, SD~2.1 y SDXL similares entre sí? Los resultados muestran que la respuesta depende del nivel de análisis. A nivel geométrico global ---dimensionalidad intrínseca, coherencia perceptual--- los tres espacios son cualitativamente equivalentes: todos son intrínsecamente de baja dimensión y todos preservan información perceptualmente relevante con correlaciones latente-imagen de entre 0,68 y 0,80. A nivel semántico, los espacios son cualitativamente distintos: el doble codificador de texto de SDXL induce una organización semántica que no se consigue simplemente escalando el UNet ni aumentando la dimensión de un único codificador.

El fine-tuning sobre WikiArt añade una dimensión inesperada a este cuadro. La organización semántica de SDXL, que lo distingue favorablemente en pretrain, resulta ser precisamente su punto de mayor vulnerabilidad ante la adaptación de dominio: una representación más estructurada puede ser más disruptible cuando el dominio nuevo es cualitativamente opuesto al del preentrenamiento. SD~1.5 y SD~2.1, cuya menor organización podría interpretarse como una limitación, exhiben en cambio una robustez geométrica y semántica casi total ante el mismo fine-tuning. La plasticidad es, paradójicamente, consecuencia de la menor rigidez representacional.

Esta asimetría ---convergencia en las propiedades geométricas globales pretrain, divergencia en las propiedades semánticas, e inversión del orden de fragilidad bajo fine-tuning--- sugiere que la calidad del condicionamiento textual es el factor determinante de la organización semántica pretrain, y que esa misma organización determina la susceptibilidad a la disrupción por adaptación de dominio. Verificar si este patrón se mantiene en modelos de la siguiente generación (DiT, flow matching) y con dominios de fine-tuning menos extremos que la pintura clásica es quizás la extensión natural más directa e informativa de este trabajo.

% ===========================================================
%  APÉNDICES
% ===========================================================
\appendix

\label{app:prompts}

\section{Prompts base}

\begin{longtable}{clp{9.5cm}}
\caption{Lista completa de los 20 prompts base del corpus (diseñados con Claude, Anthropic, 2024).}
\label{tab:prompts_completos} \\
\toprule
\textbf{ID} & \textbf{Dominio} & \textbf{Prompt} \\
\midrule
\endfirsthead
\toprule \textbf{ID} & \textbf{Dominio} & \textbf{Prompt} \\ \midrule
\endhead
\midrule \multicolumn{3}{r}{\small\textit{continúa\ldots}} \endfoot
\bottomrule \endlastfoot
base\_01 & Naturaleza muerta & \textit{a red apple on a wooden table, studio lighting, photorealistic} \\
base\_02 & Animal / exterior & \textit{a golden retriever dog sitting in a sunlit park} \\
base\_03 & Urbano / noche & \textit{a futuristic city skyline at night with neon lights} \\
base\_04 & Paisaje natural & \textit{a mountain lake at sunrise with mist and reflections} \\
base\_05 & Interior & \textit{a minimalist white living room with afternoon light} \\
base\_06 & Arquitectura & \textit{an ancient stone castle on a hilltop, overcast sky} \\
base\_07 & Tropical / costa & \textit{a tropical beach with turquoise water and palm trees} \\
base\_08 & Bosque / otoño & \textit{a forest path in autumn with golden leaves} \\
base\_09 & Bodegón bebida & \textit{a cup of coffee with latte art, warm morning light} \\
base\_10 & Tecnología / cielo & \textit{a rocket launching at dusk, long exposure photography} \\
base\_11 & Gastronomía & \textit{a bowl of ramen with steam rising, food photography} \\
base\_12 & Animal / interior & \textit{a black cat on a windowsill with rain outside} \\
base\_13 & Objetos / urbano & \textit{a vintage red bicycle against a textured brick wall} \\
base\_14 & Naturaleza / agua & \textit{a waterfall in a lush tropical jungle} \\
base\_15 & Fauna / sabana & \textit{an elephant in the African savanna at golden hour} \\
base\_16 & Arquitectura religiosa & \textit{Gothic cathedral interior with stained glass windows} \\
base\_17 & Montaña / noche & \textit{a snow-covered mountain peak under a starry night sky} \\
base\_18 & Cabaña / invierno & \textit{a wooden cabin in a snowy pine forest} \\
base\_19 & Campo / verano & \textit{a sunflower field under a vivid blue sky} \\
base\_20 & Arquitectura moderna & \textit{a modern glass skyscraper reflecting blue sky} \\
\end{longtable}

\section{Prompts de atributo}

\begin{longtable}{llp{4.4cm}p{4.4cm}}
\caption{Lista completa de los 30 prompts de atributo (15 grupos $\times$ 2 valores, diseñados con Claude, Anthropic, 2024).}
\label{tab:prompts_atributo} \\
\toprule
\textbf{Tipo} & \textbf{Grupo} & \textbf{Valor 1} & \textbf{Valor 2} \\
\midrule
\endfirsthead
\toprule \textbf{Tipo} & \textbf{Grupo} & \textbf{Valor 1} & \textbf{Valor 2} \\ \midrule
\endhead
\midrule \multicolumn{4}{r}{\small\textit{continúa\ldots}} \endfoot
\bottomrule \endlastfoot
Color & color\_apple & \textit{a warm red apple on a white surface} & \textit{a cool blue apple on a white surface} \\
Color & color\_car & \textit{a warm orange sports car on the road} & \textit{a cool blue sports car on the road} \\
Color & color\_umbrella & \textit{a warm yellow umbrella in urban rain} & \textit{a cool teal umbrella in urban rain} \\
Color & color\_flower & \textit{a warm red rose in a vase} & \textit{a cool blue rose in a vase} \\
Color & color\_balloon & \textit{a warm orange hot air balloon} & \textit{a cool blue hot air balloon} \\
Ilum. & light\_portrait & \textit{portrait, bright high-key lighting} & \textit{portrait, dark low-key lighting} \\
Ilum. & light\_kitchen & \textit{modern kitchen, bright morning light} & \textit{modern kitchen, dark night LED light} \\
Ilum. & light\_forest & \textit{forest path, bright midday sunlight} & \textit{forest path, dark sunset light} \\
Ilum. & light\_street & \textit{cobblestone street, bright sunny afternoon} & \textit{cobblestone street, dark night with lamppost} \\
Ilum. & light\_still\_life & \textit{fruit still life, bright studio lighting} & \textit{fruit still life, single candle light} \\
Estilo & style\_mountain & \textit{snowy mountain, photorealistic} & \textit{snowy mountain, impressionist painting} \\
Estilo & style\_portrait & \textit{young man portrait, photorealistic} & \textit{young man portrait, watercolor style} \\
Estilo & style\_city & \textit{aerial city view, photorealistic} & \textit{aerial city view, digital art style} \\
Estilo & style\_forest & \textit{light through trees, photorealistic} & \textit{light through trees, fantasy painting} \\
Estilo & style\_cat & \textit{cat on sofa, photorealistic} & \textit{cat on sofa, Studio Ghibli style} \\
\end{longtable}

% ===========================================================
%  REFERENCIAS BIBLIOGRÁFICAS  (APA 7)
% ===========================================================
\chapter*{Referencias bibliográficas}
\addcontentsline{toc}{chapter}{Referencias bibliográficas}

\begin{itemize}[leftmargin=1.5em, itemsep=4pt]

\item Anthropic. (2024). \textit{Claude} [Modelo de lenguaje de gran escala]. Anthropic. \url{https://www.anthropic.com}

\item Betker, J., Goh, G., Jing, L., Brooks, T., Wang, J., Li, L., \ldots{} Ramesh, A. (2023). \textit{Improving image generation with better captions} [Informe técnico]. OpenAI.

\item Black Forest Labs. (2024). \textit{Flux.1: A family of flow matching text-to-image models}. \url{https://github.com/black-forest-labs/flux}

\item Dinh, L., Sohl-Dickstein, J., y Bengio, S. (2017). Density estimation using real-valued non-volume preserving transformations. \textit{International Conference on Learning Representations}.

\item Efros, A. A., y Leung, T. K. (1999). Texture synthesis by non-parametric sampling. \textit{Proceedings of the IEEE International Conference on Computer Vision}, 1033--1038.

\item Esser, P., Kulal, S., Blattmann, A., Entezari, R., M{\"u}ller, J., Saini, H., \ldots{} Rombach, R. (2024). Scaling rectified flow transformers for high-resolution image synthesis. \textit{Proceedings of the 41st International Conference on Machine Learning}.

\item Facco, E., d'Errico, M., Rodriguez, A., y Laio, A. (2017). Estimating the intrinsic dimension of datasets by a minimal neighborhood information. \textit{Scientific Reports, 7}(1), Art.~12140. \url{https://doi.org/10.1038/s41598-017-11873-y}

\item Gal, R., Alaluf, Y., Atzmon, Y., Patashnik, O., Bermano, A. H., Chechik, G., y Cohen-Or, D. (2022). An image is worth one word: Personalizing text-to-image generation using textual inversion. \textit{International Conference on Learning Representations 2023}.

\item Gatys, L. A., Ecker, A. S., y Bethge, M. (2016). Image style transfer using convolutional neural networks. \textit{IEEE/CVF Conference on Computer Vision and Pattern Recognition}, 2414--2423.

\item Google LLC. (2017). \textit{Google Colaboratory}. \url{https://colab.research.google.com}

\item Goodfellow, I., Pouget-Abadie, J., Mirza, M., Xu, B., Warde-Farley, D., Ozair, S., Courville, A., y Bengio, Y. (2014). Generative adversarial nets. \textit{Advances in Neural Information Processing Systems, 27}.

\item H{\"a}rkonen, E., Hertzmann, A., Lehtinen, J., y Paris, S. (2020). GANSpace: Discovering interpretable GAN controls. \textit{Advances in Neural Information Processing Systems, 33}.

\item Hertz, A., Mokady, R., Tenenbaum, J., Aberman, K., Pritch, Y., y Cohen-Or, D. (2022). Prompt-to-Prompt image editing with cross attention control. \textit{International Conference on Learning Representations 2023}.

\item Ho, J., Jain, A., y Abbeel, P. (2020). Denoising diffusion probabilistic models. \textit{Advances in Neural Information Processing Systems, 33}, 6840--6851.

\item Ho, J., y Salimans, T. (2022). Classifier-free diffusion guidance. \textit{NeurIPS 2021 Workshop on Deep Generative Models and Downstream Applications}.

\item Hu, E. J., Shen, Y., Wallis, P., Allen-Zhu, Z., Li, Y., Wang, S., Wang, L., y Chen, W. (2022). LoRA: Low-rank adaptation of large language models. \textit{International Conference on Learning Representations}.

\item Ilharco, G., Wortsman, M., Wightman, R., Gordon, C., Carlini, N., Taori, R., \ldots{} Schmidt, L. (2021). \textit{OpenCLIP}. Zenodo. \url{https://doi.org/10.5281/zenodo.5143773}

\item Karras, T., Laine, S., y Aila, T. (2019). A style-based generator architecture for generative adversarial networks. \textit{IEEE/CVF Conference on Computer Vision and Pattern Recognition}, 4401--4410.

\item Karras, T., Laine, S., Aittala, M., Hellsten, J., Lehtinen, J., y Aila, T. (2020). Analyzing and improving the image quality of StyleGAN. \textit{IEEE/CVF Conference on Computer Vision and Pattern Recognition}, 8110--8119.

\item Kingma, D. P., y Dhariwal, P. (2018). Glow: Generative flow with invertible 1$\times$1 convolutions. \textit{Advances in Neural Information Processing Systems, 31}.

\item Kingma, D. P., y Welling, M. (2013). Auto-encoding variational Bayes. \textit{International Conference on Learning Representations 2014}.

\item Kirkpatrick, J., Pascanu, R., Rabinowitz, N., Veness, J., Desjardins, G., Rusu, A. A., \ldots{} Hadsell, R. (2017). Overcoming catastrophic forgetting in neural networks. \textit{Proceedings of the National Academy of Sciences, 114}(13), 3521--3526.

\item Krizhevsky, A., Sutskever, I., y Hinton, G. E. (2012). ImageNet classification with deep convolutional neural networks. \textit{Advances in Neural Information Processing Systems, 25}.

\item Kwon, M., Jeong, J., y Uh, Y. (2023). Diffusion models already have a semantic latent space. \textit{International Conference on Learning Representations}.

\item Lipman, Y., Chen, R. T. Q., Ben-Hamu, H., Nickel, M., y Le, M. (2022). Flow matching for generative modeling. \textit{International Conference on Learning Representations 2023}.

\item Liu, X., Gong, C., y Liu, Q. (2022). Flow straight and fast: Learning to generate and transfer data with rectified flow. \textit{International Conference on Learning Representations 2023}.

\item McCloskey, M., y Cohen, N. J. (1989). Catastrophic interference in connectionist networks: The sequential learning problem. \textit{Psychology of Learning and Motivation, 24}, 109--165.

\item Mokady, R., Hertz, A., Aberman, K., Pritch, Y., y Cohen-Or, D. (2023). Null-text inversion for editing real images using guided diffusion models. \textit{IEEE/CVF Conference on Computer Vision and Pattern Recognition}.

\item Peebles, W., y Xie, S. (2023). Scalable diffusion models with transformers. \textit{IEEE/CVF International Conference on Computer Vision}.

\item Perlin, K. (1985). An image synthesizer. \textit{ACM SIGGRAPH Computer Graphics, 19}(3), 287--296.

\item Podell, D., English, Z., Lacey, K., Blattmann, A., Dockhorn, T., M{\"u}ller, J., Penna, J., y Rombach, R. (2023). SDXL: Improving latent diffusion models for high-resolution image synthesis. \textit{arXiv}. \url{https://arxiv.org/abs/2307.01952}

\item Pope, P., Zhu, C., Abdelkader, A., Goldblum, M., y Goldstein, T. (2021). The intrinsic dimension of images and its impact on learning. \textit{International Conference on Learning Representations}.

\item Radford, A., Metz, L., y Chintala, S. (2015). Unsupervised representation learning with deep convolutional generative adversarial networks. \textit{arXiv}. \url{https://arxiv.org/abs/1511.06434}

\item Radford, A., Kim, J. W., Hallacy, C., Ramesh, A., Goh, G., Agarwal, S., \ldots{} Sutskever, I. (2021). Learning transferable visual models from natural language supervision. \textit{International Conference on Machine Learning}, 8748--8763.

\item Rombach, R., Blattmann, A., Lorenz, D., Esser, P., y Ommer, B. (2022). High-resolution image synthesis with latent diffusion models. \textit{IEEE/CVF Conference on Computer Vision and Pattern Recognition}, 10684--10695.

\item Ruiz, N., Li, Y., Jampani, V., Pritch, Y., Rubinstein, M., y Aberman, K. (2023). DreamBooth: Fine tuning text-to-image diffusion models for subject-driven generation. \textit{IEEE/CVF Conference on Computer Vision and Pattern Recognition}.

\item Saleh, B., y Elgammal, A. (2015). Large-scale classification of fine-art paintings: Learning the right metric on the right feature. \textit{arXiv}. \url{https://arxiv.org/abs/1505.00855}

\item Schuhmann, C., Beaumont, R., Vencu, R., Gordon, C., Wightman, R., Cherti, M., \ldots{} Jitsev, J. (2022). LAION-5B: An open large-scale dataset for training next generation image-text models. \textit{Advances in Neural Information Processing Systems, 35}.

\item Shen, Y., Gu, J., Tang, X., y Zhou, B. (2020). Interpreting the latent space of GANs for semantic face editing. \textit{IEEE/CVF Conference on Computer Vision and Pattern Recognition}, 9243--9252.

\item Shen, Y., y Zhou, B. (2021). Closed-form factorization of latent semantics in GANs. \textit{IEEE/CVF Conference on Computer Vision and Pattern Recognition}.

\item Song, J., Meng, C., y Ermon, S. (2021). Denoising diffusion implicit models. \textit{International Conference on Learning Representations}.

\item Stability AI. (2022). \textit{Stable Diffusion 2}. \url{https://github.com/Stability-AI/stablediffusion}

\item Tang, L., Jia, M., Wang, Q., Phoo, C. P., y Hariharan, B. (2023). Emergent correspondence from image diffusion. \textit{Advances in Neural Information Processing Systems, 36}.

\item van den Oord, A., Kalchbrenner, N., y Kavukcuoglu, K. (2016). Pixel recurrent neural networks. \textit{Proceedings of the 33rd International Conference on Machine Learning}.

\item Wang, Z., Bovik, A. C., Sheikh, H. R., y Simoncelli, E. P. (2004). Image quality assessment: From error visibility to structural similarity. \textit{IEEE Transactions on Image Processing, 13}(4), 600--612.

\item Wu, Z., Lischinski, D., y Shechtman, E. (2021). StyleSpace analysis: Disentangled controls for StyleGAN image generation. \textit{IEEE/CVF Conference on Computer Vision and Pattern Recognition}.

\end{itemize}

\end{document}
